% -----------------------------------------------------------
\section{1836}
% -----------------------------------------------------------
	\label{EMS-1836}
	
	% ----------------------
	\subsection{Joseph Smith}
	% ----------------------
		\label{EMS-1836-JS}

		% -----
		\subsubsection{Introduction to 1836}
		% -----
		
		% -----
		\subsubsection{Record, 1 January 1836}
		% -----
			\label{EMS-1836-JS-1-JAN-1836}
			% \label{EMS-1830-JS-DAY-3LETTERMONTH-YEAR}
			
			\begin{description}
					\item[Print Sources]
				\end{description}

					\begin{tabular}{p{0.5in}p{1in}p{0.5in}p{1in}}
					JSP	 	& J1:140--141		& JSR 		& --- \\
					DC	 	& ---				& HC 		& --- \\
					HDDC 	& ---				& TCHC 		& --- \\
					PJS 	& 1:181--183			& 	 		&  \\
					\end{tabular}
					% HDDC = woodford's DC diss; JSR = Marquardt's JS Revelations; TCHC = Vogel HC
				
				\begin{description}
					\item[Online Access] \href{https://www.josephsmithpapers.org/paper-summary/journal-1835-1836/95}{Here}
					% \item[Analysis] \hyperref[XX]{Here}
				\end{description}
				
				\begin{description}
					\item[Background]
				\end{description}
				
					% It might also be useful to give a two sentence context of the period: e.g., "This speech was given in Nauvoo to a small congregation one month prior to the destruction of the Nauvoo Expositor. These and other tensions may have been on the prophet's mind when he gave the speech."
				
				\begin{description}
					\item[Original Source]
				\end{description}
				
					% brief description of original source material, with reference to JSP or somewhere else for more info
					% Background material: it might be helpful to have a note that says whether a given document was presented as a speech, was written in a journal, etc.
					% Also a comment here about how reliable the source might be, or what shape it is in, if there are large omissions or parts that are hard to understand, and so on.
				
				\begin{description}
					\item[Text]
				\end{description}
				
						\hypertarget{EMS-1836-JS-1-JAN-1836-a}%
					Friday
						\marginnote{a}%
					morning Jan\supers{y.} 1\supers{st.} 1836 this being the beginning of a new year, my heart is filled with gratitude to God, that he has preserved my life and the \sout{life} <lives> of my family while another year has rolled away, we have been, sustained and upheld in the midst of a wicked and perverse generation, and exposed to all, the afflictions temptations and misery that are incident to human life, for which I feel to humble myself in dust and ashes, as it were before the Lord--- but notwithstanding, the gratitude that fills my heart on retrospecting the past year, and the multiplyed blessings that have crowned our heads, my heart is pained within me because of the difficulty that \sout{in} exists in my fathers family, the Devil has made a violent attack on Br. W\supersu{m} [Smith] and Br Calvin [Stoddard] and the powers of darkness, seeme [to] lower over their minds and not only theirs but cast a gloomy shade over the minds of my \sout{my} \sout{parents} \sout{and} \sout{somee} \sout{of} \sout{my} brothers and sisters, which prevents them from seeing things as they realy are, and the powers of Earth \& hell seem combined to overthrow us and the Church by causing a division in the family, and \sout{is} indeed the adversary is bring[ing] into requisition all his subtlety to prevent the Saints from being endowed, by causing devision among the 12, also among the 70, and bickerings and jealousies among the Elders and official members of the church, and so the leaven of iniquity foments and spreads among the members of the church,
					
						\hypertarget{EMS-1836-JS-1-JAN-1836-b}%
					But
						\marginnote{b}%
					I am determined that nothing on my part shall be lacking to adjust and amicably dispose of and settle all family difficulties, on this day, that the ensuing year, and years, be they many or few may be spent in righteousneess before God, and I know that the cloud will burst and satans kingdom be laid in \sout{ruin,} <ruins> with all his black designs, and the saints come forth like gold seven times tried in the fire, being made perfect throug[h] sufferings, and temptations, and the blessings of heaven and earth multiplyed upon our heads which may God grant for Christ sake Amen---
					
						\hypertarget{EMS-1836-JS-1-JAN-1836-c}%
					Br.
						\marginnote{c}%
					William came to my house and Br. Hyrum [Smith], also, Uncle John Smith, we went into a room in company with father and Elder Martin Harris, \sout{and} father, Smith then opened our interview by prayer after which, he expressed his feelings on the ocasion in a verry feeling and pathetic manner even with all the sympathy of a father whose feeling were wounded deeply on the account of the difficulty that was existing in the family, and while he addressed us the spirit of God rested down upon us in mighty power, and our hearts were melted Br. William made an humble confession and asked \sout{our} my forgiveness for the abuse he had offered me and wherein I had been out of the way I asked his forgivness, and the spirit of confission and forgiveness, was mutual among us all, and we covenanted with each other in the Sight of God and the holy angels and the brethren, to strive from henceforward to build each other up in righteousness, in all things and not listen to evil reports concerning eachother, but like brethren, indeed go to eachother, with our grievances in the spirit of meekness, and be reconciled and thereby promote our own happiness and the happiness of the family and in short the happiness and well being of all.--- my wife and Mother, \sout{Uncle} \sout{John} \& my Scribe was then called in and we repeated the covenant to them that we had entered into, and while gratitude swelled our bosoms, tears flowed from our eys.--- I was then requested to close our interview which I did with prayer, and it was truly a jubilee and time of rejoiceing
	
		% -----
		\subsubsection{Blessing to Lorenzo Barnes, 3 January 1836}
		% -----
			\label{EMS-1836-JS-3-JAN-1836}
			% \label{EMS-1830-JS-DAY-3LETTERMONTH-YEAR}
			
			\begin{description}
					\item[Print Sources]
				\end{description}

					\begin{tabular}{p{0.5in}p{1in}p{0.5in}p{1in}}
					JSP	 	& D5:131--135		& JSR 		& --- \\
					DC	 	& ---				& HC 		& --- \\
					HDDC 	& ---				& TCHC 		& --- \\
					\end{tabular}
					% HDDC = woodford's DC diss; JSR = Marquardt's JS Revelations; TCHC = Vogel HC
				
				\begin{description}
					\item[Online Access] \href{https://www.josephsmithpapers.org/paper-summary/blessing-to-lorenzo-barnes-3-january-1836/1}{Here}
					% \item[Analysis] \hyperref[XX]{Here}
				\end{description}
				
				\begin{description}
					\item[Background]
				\end{description}
				
					% It might also be useful to give a two sentence context of the period: e.g., "This speech was given in Nauvoo to a small congregation one month prior to the destruction of the Nauvoo Expositor. These and other tensions may have been on the prophet's mind when he gave the speech."
				
				\begin{description}
					\item[Original Source]
				\end{description}
				
					% brief description of original source material, with reference to JSP or somewhere else for more info
					% Background material: it might be helpful to have a note that says whether a given document was presented as a speech, was written in a journal, etc.
					% Also a comment here about how reliable the source might be, or what shape it is in, if there are large omissions or parts that are hard to understand, and so on.
				
				\begin{description}
					\item[Text]
				\end{description}
				
						\hypertarget{EMS-1836-JS-3-JAN-1836-a}%
					L[orenzo]
						\marginnote{a}%
					Barnes, Zion Blessing received under the hands of Joseph Smith J\supersu{r}\supers{.} Sidney Rigdon \& Hyram [Hyrum] Smith Presidents of the church of Latter Day Saints Kirtland Ohio J\supersu{an}\supers{.} 3\supersu{d} 1836
					
					Brother Barnes
					
					we lay our hands <up>on thy head in the name of the Lord \& in his name we say unto thee thou art excepted [accepted] before him The Eyes of the Lord thy God hav been upon thee \& thou hast done that which was most pleasing in his sight for thou hast in thy youth set out in his serves [service] \sout{\&} Thy Prayers \& suplications hav been herd \& thy Name is writen in Heaven for Angels to gaze upon and thou shalt be a swift messinger to the Nations.
					
						\hypertarget{EMS-1836-JS-3-JAN-1836-b}%
					Thou
						\marginnote{b}%
					art a desendant of Joseph \& of the Tribe of Ep[h]raim \& the blessings of Jacob \& Joseph are thine even the choise Blessings of Heaven above \& the Earth beneath \& the fullness there of Because thou hast been faithful and hast not withheld thy Life from layinging it down for thy breathren Thou art a chosen vessel unto the Lord to bare the fullness of the Gosple unto people \& Nations a far off Be faithful \& thou shalt be endowed with power from on high for the spirit of the highest shall rest upon thee \& thou shalt go forth from Land to Land from Nation to Nation \& from Kingdom to Kingdom Thou shalt Stand before Kings \& \uline{\sout{rulers}} <\uline{Princes}> the rich \& the great \& they shall be estonished at thy wisdom \& tremble \sout{at} at thy words for the Lord thy God will unloos thy tongue hitherto thou hast been week but thou shalt be made strong Thou shalt be mighty in speaking Thy voice shall bee like the voice of an Angel of God Thy faith shall be mighty in the Earth Thou shalt lift up thy voice \& the hills shall tremble before the[e] \& the Earth shall quake under thy feet Prison walls shall not hold thee watters shall not stay thee The sea \& the wind shall obey thy command Thou shalt lay thine hands on the sick \& they shall recover Thou shalt hav power to open the eyes of the blind \& unstop the ears of the deaf Thou shalt cause the lame to leap as an heart [hart] \& the tongue of the dumb to sing yea thou shalt bring thy thousands \& thy \& thy tens of thousands from the Nations to Zion with rejoicing \& sit down with them \sout{in} <\&> with Abraham Isaac \& Jacob in the Kingdom of God
					
						\hypertarget{EMS-1836-JS-3-JAN-1836-c}%
					Therefore
						\marginnote{c}%
					be faithful \& thy God shall lift the[e] up \& thou shalt be exalted \& he shall say of <thee> This is my servant <in my sight> he is highley exalted \sout{in my sight} Therefore we say unto thee our brother let thy mind expand for it is thy privaleg to behold the Heavens opened over thy head to see \& converce with the Angel of God for they shall appear unto thee in thy ministry y[e]a thou shalt behold the fase of thy Redeemer in the flesh Thou shalt hav communion with the general assembly \& Church of the first born in Heaven y[e]a thou shalt behold with thine eyes the glory of \sout{thy} the Lord untill thou shhalt be satisfied \& thou shalt say it is \sout{a} enough O Lord it is enough Thou shalt hav all the desires of thy heart Thou shalt hav length of days \& thy years shall be many
					
						\hypertarget{EMS-1836-JS-3-JAN-1836-d}%
					When
						\marginnote{d}%
					thy hears [hairs] shall be gray thou shalt be as in the vigur of youth \& thy mind unimpared Thou shalt remane untill the coming of the Son of man in the clowds of heaven y[e]a thou shalt see <view> the winding up seen of all things \& stand with the hundred forty \& four thousand on Mount Zion These blessings \& many others which tongue cannot express \& all the powers of the priesthood we seal upon thy head \& upon thy sead forever in the name of Jesus Christ \& by the \sout{power} arthority of the High Priesthood even so amen
					
					A\supersu{pr}\supers{.} 24\supersu{th} 1836 [\textit{1/3 page blank}]
					
		% -----
		\subsubsection{Record, 7 January 1836}
		% -----
			\label{EMS-1836-JS-7-JAN-1836}
			% \label{EMS-1830-JS-DAY-3LETTERMONTH-YEAR}
			
			\begin{description}
					\item[Print Sources]
				\end{description}

					\begin{tabular}{p{0.5in}p{1in}p{0.5in}p{1in}}
					JSP	 	& J1:146		& JSR 		& --- \\
					DC	 	& ---				& HC 		& --- \\
					HDDC 	& ---				& TCHC 		& --- \\
					PJS 	& 1:186--187			& 	 		&  \\
					\end{tabular}
					% HDDC = woodford's DC diss; JSR = Marquardt's JS Revelations; TCHC = Vogel HC
				
				\begin{description}
					\item[Online Access] \href{https://www.josephsmithpapers.org/paper-summary/journal-1835-1836/102}{Here}
					% \item[Analysis] \hyperref[XX]{Here}
				\end{description}
				
				\begin{description}
					\item[Background]
				\end{description}
				
					% It might also be useful to give a two sentence context of the period: e.g., "This speech was given in Nauvoo to a small congregation one month prior to the destruction of the Nauvoo Expositor. These and other tensions may have been on the prophet's mind when he gave the speech."
				
				\begin{description}
					\item[Original Source]
				\end{description}
				
					% brief description of original source material, with reference to JSP or somewhere else for more info
					% Background material: it might be helpful to have a note that says whether a given document was presented as a speech, was written in a journal, etc.
					% Also a comment here about how reliable the source might be, or what shape it is in, if there are large omissions or parts that are hard to understand, and so on.
				
				\begin{description}
					\item[Text]
				\end{description}
				
					Thursday 7\supers{th} attended a sumptuous feast at Bishop N[ewel] K. Whitneys this feast was after the order of the Son of God the lame the halt and blind wer invited according to the intruction of the Saviour
					
					our meeting was opened by singing and prayer offered up by father Smith, after which Bishop, Whitneys father \& mother were bless[ed] and a number of others, with a patriarchal blessing, we then recieved a bountiful refreshment, furnished by the liberality of the Bishop the company was large,--- before we parted we had some of the Songs of Zion sung, and our hearts were made glad while partaking of an antipast of those Joys that will be poured upon the head of the Saints w[h]en they are gathered together on Mount Zion to enjoy eachothers society forever more even all the blessings of heaven and earth \sout{and} where there will be none to molest nor make us afraid...
					
		% -----
		\subsubsection{Record, 12 January 1836}
		% -----
			\label{EMS-1836-JS-12-JAN-1836}
			% \label{EMS-1830-JS-DAY-3LETTERMONTH-YEAR}
			
			\begin{description}
					\item[Print Sources]
				\end{description}

					\begin{tabular}{p{0.5in}p{1in}p{0.5in}p{1in}}
					JSP	 	& J1:147--148		& JSR 		& --- \\
					DC	 	& ---				& HC 		& --- \\
					HDDC 	& ---				& TCHC 		& --- \\
					PJS 	& 1:188--189			& 	 		&  \\
					\end{tabular}
					% HDDC = woodford's DC diss; JSR = Marquardt's JS Revelations; TCHC = Vogel HC
				
				\begin{description}
					\item[Online Access] \href{https://www.josephsmithpapers.org/paper-summary/journal-1835-1836/105}{Here}
					% \item[Analysis] \hyperref[XX]{Here}
				\end{description}
				
				\begin{description}
					\item[Background]
				\end{description}
				
					% It might also be useful to give a two sentence context of the period: e.g., "This speech was given in Nauvoo to a small congregation one month prior to the destruction of the Nauvoo Expositor. These and other tensions may have been on the prophet's mind when he gave the speech."
				
				\begin{description}
					\item[Original Source]
				\end{description}
				
					% brief description of original source material, with reference to JSP or somewhere else for more info
					% Background material: it might be helpful to have a note that says whether a given document was presented as a speech, was written in a journal, etc.
					% Also a comment here about how reliable the source might be, or what shape it is in, if there are large omissions or parts that are hard to understand, and so on.
				
				\begin{description}
					\item[Text]
				\end{description}
				
					also a man was introduced to me by the name of Russel Wever [Russell Weaver] from Cambray Niagary [Cambria, Niagara] Co. N. Y this man is a preacher, in the church that is called Christian or Unitarian, \sout{some} he remarked that he had but few minuits to spend with me, we entered into conversation, and had some little controversy upon the subject of prejudice, but soon come to \sout{the} an understanding, he spoke of the gospel and said he believed it, adding that it was good tidings of great joy--- I replyed that it was one thing, to proclaim good tidings and another to tell what those tidings are, he waived the conversation and withdrew--- he was introduced by Joseph Rose---
					
		% -----
		\subsubsection{Rules and Regulations, 14 January 1836}
		% -----
			\label{EMS-1836-JS-14-JAN-1836}
			% \label{EMS-1830-JS-DAY-3LETTERMONTH-YEAR}
			
			\begin{description}
					\item[Print Sources]
				\end{description}

					\begin{tabular}{p{0.5in}p{1in}p{0.5in}p{1in}}
					JSP	 	& D5:143--145; J1:152--153		& JSR 		& --- \\
					DC	 	& ---				& HC 		& 2:368--369 \\
					HDDC 	& ---				& TCHC 		& 2:359 \\
					PJS 	& 1:193--197			& 	 		&  \\
					\end{tabular}
					% HDDC = woodford's DC diss; JSR = Marquardt's JS Revelations; TCHC = Vogel HC
				
				\begin{description}
					\item[Online Access] \href{https://www.josephsmithpapers.org/paper-summary/rules-and-regulations-14-january-1836/1}{Here}
					% \item[Analysis] \hyperref[XX]{Here}
				\end{description}
				
				\begin{description}
					\item[Background]
				\end{description}
				
					See the J1 entry for 13 January, the day before (J1:150).
				
					% It might also be useful to give a two sentence context of the period: e.g., "This speech was given in Nauvoo to a small congregation one month prior to the destruction of the Nauvoo Expositor. These and other tensions may have been on the prophet's mind when he gave the speech."
				
				\begin{description}
					\item[Original Source]
				\end{description}
				
					% brief description of original source material, with reference to JSP or somewhere else for more info
					% Background material: it might be helpful to have a note that says whether a given document was presented as a speech, was written in a journal, etc.
					% Also a comment here about how reliable the source might be, or what shape it is in, if there are large omissions or parts that are hard to understand, and so on.
				
				\begin{description}
					\item[Text]
				\end{description}
				
						\hypertarget{EMS-1836-JS-14-JAN-1836-a}%
					1\supers{st}---
						\marginnote{a}%
					It is according to the rules and regulations of all regular and legal organized bodies to have a president to keep order.---

					2\supers{ond}--- The body thus organized are under obligation to be in subjection to that authority---

					3\supers{d}--- When a congregation assembles in this house they shall submit to the following rules, that due respect may be paid to the order of the worship---viz.
					
					\phantom{.}
					
						\hypertarget{EMS-1836-JS-14-JAN-1836-b}%
					1\supers{st}---
						\marginnote{b}%
					no man shall be interupted who is appointed to speak by the presidency of the Church, by any disorderly person or persons in the congregation, by whispering by laughing by talking by menacing Jestures by getting up and running out in a disorderly manner or by offering indignity to the manner of worship or the religion or to any officer of said church while officiating in his office in any wise whatever by any display of ill manners or ill breeding from old or young rich or poor male or female bond or free black or white believer or unbeliever and if any of the above insults are offered such measures will be taken as are lawful to punish the aggressor or aggressors and eject them out of the house

					2\supers{ond}--- An insult offered to the \sout{presidency} <presiding> Elder of said Church, shall be concidered an insult to the whole body, also an insult offered to any of the officers of said Church while officiating shall be considered an insult to the whole body---

					3\supers{d}--- All persons are prohibited from going up the stairs in times of worship

					4\supers{th}--- all persons are prohibited from exploring the house except waited upon by a person appointed for that purpose---

						\hypertarget{EMS-1836-JS-14-JAN-1836-c}%
					5\supers{th}---
						\marginnote{c}%
					all persons are prohibited from going <in>to the several pulpits except the officers who are appointed to officiate in the same

					6\supers{th}--- All persons are prohibited from cutting marking or marring the inside or outside of the house with a knife pencil or any other instrument whatever, under pain of such penalty as the law shall inflict---

					7\supers{th}--- All children are prohibited from assembling in the house above or below or any part of it to play or for recreation at any time, and all parents guardians or masters shall be ameneable for all damage that shall accrue in consequence of their children---

					8\supers{th}--- All persons whether believers or unbelievers shall be treated with due respect by the \sout{authority} <authorities> of the Church---

					9\supers{th}--- no imposition shall be practiced upon any member of the church by depriving them of their <rights> in the house--- council adjourned sini di [sine die]
					
		% -----
		\subsubsection{Record, 14 January 1836}
		% -----
			\label{EMS-1836-JS-14-JAN-1836-REC}
			% \label{EMS-1830-JS-DAY-3LETTERMONTH-YEAR}
			
			\begin{description}
					\item[Print Sources]
				\end{description}

					\begin{tabular}{p{0.5in}p{1in}p{0.5in}p{1in}}
					JSP	 	& J1:153		& JSR 		& --- \\
					DC	 	& ---				& HC 		& --- \\
					HDDC 	& ---				& TCHC 		& --- \\
					PJS 	& 1:197			& 	 		&  \\
					\end{tabular}
					% HDDC = woodford's DC diss; JSR = Marquardt's JS Revelations; TCHC = Vogel HC
				
				\begin{description}
					\item[Online Access] \href{https://www.josephsmithpapers.org/paper-summary/journal-1835-1836/114}{Here}
					% \item[Analysis] \hyperref[XX]{Here}
				\end{description}
				
				\begin{description}
					\item[Background]
				\end{description}
				
					% It might also be useful to give a two sentence context of the period: e.g., "This speech was given in Nauvoo to a small congregation one month prior to the destruction of the Nauvoo Expositor. These and other tensions may have been on the prophet's mind when he gave the speech."
				
				\begin{description}
					\item[Original Source]
				\end{description}
				
					% brief description of original source material, with reference to JSP or somewhere else for more info
					% Background material: it might be helpful to have a note that says whether a given document was presented as a speech, was written in a journal, etc.
					% Also a comment here about how reliable the source might be, or what shape it is in, if there are large omissions or parts that are hard to understand, and so on.
				
				\begin{description}
					\item[Text]
				\end{description}
				
					...being under obligation to attend at Mrs. Wilcoxs to join Mr. John Webb and Mrs Catharine [Catherine] Wilcox in matrimony also M\supers{r.} Th\supers{os} Carrico [Jr.] and Miss Elizabeth Baker at the same place, I found a large company assembled, the house was filled to overflowing, we opened our interview by singing and prayer suited to the occasion after which I made some remarks in relation to the duties that are incumbent on husbands and wives, in particular the great importance there is in cultivating the pure principles of the institution, in all it's bearings, and connexions with each other and society in general
					
					I then invited them to arise and join hands, and pronounced the ceremony according to the rules and regulations of the Church of the latterday Saints
					
					\sout{Closed} after which I pronounced such blessings upon their heads as the Lord put into my heart \sout{even} even the blessings of Abraham Isaac and Jacob, and dismissed by singing and prayer...
					
		% -----
		\subsubsection{Minutes, 15 January 1836}
		% -----
			\label{EMS-1836-JS-15-JAN-1836}
			% \label{EMS-1830-JS-DAY-3LETTERMONTH-YEAR}
			
			\begin{description}
					\item[Print Sources]
				\end{description}

					\begin{tabular}{p{0.5in}p{1in}p{0.5in}p{1in}}
					JSP	 	& D5:146--148		& JSR 		& --- \\
					DC	 	& ---				& HC 		& 2:370--371 \\
					HDDC 	& ---				& TCHC 		& 2:360--363 \\
					\end{tabular}
					% HDDC = woodford's DC diss; JSR = Marquardt's JS Revelations; TCHC = Vogel HC
				
				\begin{description}
					\item[Online Access] \href{https://www.josephsmithpapers.org/paper-summary/minutes-15-january-1836/1}{Here}
					% \item[Analysis] \hyperref[XX]{Here}
				\end{description}
				
				\begin{description}
					\item[Background]
				\end{description}
				
					% It might also be useful to give a two sentence context of the period: e.g., "This speech was given in Nauvoo to a small congregation one month prior to the destruction of the Nauvoo Expositor. These and other tensions may have been on the prophet's mind when he gave the speech."
				
				\begin{description}
					\item[Original Source]
				\end{description}
				
					% brief description of original source material, with reference to JSP or somewhere else for more info
					% Background material: it might be helpful to have a note that says whether a given document was presented as a speech, was written in a journal, etc.
					% Also a comment here about how reliable the source might be, or what shape it is in, if there are large omissions or parts that are hard to understand, and so on.
				
				\begin{description}
					\item[Text]
				\end{description}
				
					Friday Morning January 15, 1836, council met pursuant to adjournment and after President J. Smith Jun\supersu{r} had organized the council, he proceeded to give many good instructions in relation to the order \& manner of conducting the council and also delivered a solemn charge to the counsel after which he opened by prayer and presided as before.
					
					President J. Smith Jun\supersu{r} one of the committee to draft rules for the regulation of the House of the Lord, made the report of said committee by reading the laws or rules they had drafted three times. Th[e]y were approved and unanimously adopted...
					
					...Motioned, seconded \& voted that the Presidencey of the high council hold the keys of the house of the Lord except the keys of one vestry which shall be held by the Bishop of Kirtland
					
					Motioned and seconded that the Laws regulating the house of the Lord go into effect from this time, and that Elder John Corril[l] take it upon him to see that they are enforced, giving him the privilege of calling as many as he choose to assist him...
					
		% -----
		\subsubsection{Record, 15 January 1836}
		% -----
			\label{EMS-1836-JS-15-JAN-1836-REC}
			% \label{EMS-1830-JS-DAY-3LETTERMONTH-YEAR}
			
			\begin{description}
					\item[Print Sources]
				\end{description}

					\begin{tabular}{p{0.5in}p{1in}p{0.5in}p{1in}}
					JSP	 	& D5:153--154		& JSR 		& --- \\
					DC	 	& ---				& HC 		& --- \\
					HDDC 	& ---				& TCHC 		& --- \\
					PJS 	& 1:198--200		& TPJS 		& 93--94 \\
					\end{tabular}
					% HDDC = woodford's DC diss; JSR = Marquardt's JS Revelations; TCHC = Vogel HC
				
				\begin{description}
					\item[Online Access] \href{https://www.josephsmithpapers.org/paper-summary/journal-1835-1836/115}{Here}
					% \item[Analysis] \hyperref[XX]{Here}
				\end{description}
				
				\begin{description}
					\item[Background]
				\end{description}
				
					% It might also be useful to give a two sentence context of the period: e.g., "This speech was given in Nauvoo to a small congregation one month prior to the destruction of the Nauvoo Expositor. These and other tensions may have been on the prophet's mind when he gave the speech."
				
				\begin{description}
					\item[Original Source]
				\end{description}
				
					% brief description of original source material, with reference to JSP or somewhere else for more info
					% Background material: it might be helpful to have a note that says whether a given document was presented as a speech, was written in a journal, etc.
					% Also a comment here about how reliable the source might be, or what shape it is in, if there are large omissions or parts that are hard to understand, and so on.
				
				\begin{description}
					\item[Text]
				\end{description}
				
						\hypertarget{EMS-1836-JS-15-JAN-1836-REC-a}%
					Friday
						\marginnote{a}%
					the 15\supers{th} at 9 oclock A.M meet in council agreeably to the adjournment, at the council room in the Chapel organized the authorities of the church agreeably to their respective offices in the same, I then made some observation respecting the order of the day, and the great responsibility we are under to transact all our buisness, in righteousness before God, inasmuch as our desisions will have a bearing upon all mankind and upon all generations to come...
					
						\hypertarget{EMS-1836-JS-15-JAN-1836-REC-b}%
					...The
						\marginnote{b}%
					above rules hav now passed through the several quorums, in their order, and passed by the unanimous vote of the whole, and are therefore received and established as a law to govern the house of the Lord in this place,--- In the investigated of this subject, I found that many who had deliberated upon this subject were darkened in their minds, which drewforth, some remarks from me, respecting the privileges of the authorities of the church, that they should, each speak in his turn, and in his place, and in his time and season, that their may be perfect order in all things, and that every man, before he, makes an objection to any, item, that is thrown before them for their concideration, should be sure that they can throw light upon the subject rather than spread darkness, and that his objections be founded in righteousness which may be done by applying ourselves closely to study the, mind and will of the Lord, whose Spirit always makes manifest, and demonstrates to the understanding of all who are in possession, of his Spirit---...
					
		% -----
		\subsubsection{Minutes, 16 January 1836}
		% -----
			\label{EMS-1836-JS-16-JAN-1836}
			% \label{EMS-1830-JS-DAY-3LETTERMONTH-YEAR}
			
			\begin{description}
					\item[Print Sources]
				\end{description}

					\begin{tabular}{p{0.5in}p{1in}p{0.5in}p{1in}}
					JSP	 	& D5:148--154; J1:156--160		& JSR 		& --- \\
					DC	 	& ---				& HC 		& 2:372--375 \\
					HDDC 	& ---				& TCHC 		& 2:365--368 \\
					TPJS 	& 105--106			& PJS	 	& 1:201--206 \\
					\end{tabular}
					% HDDC = woodford's DC diss; JSR = Marquardt's JS Revelations; TCHC = Vogel HC
				
				\begin{description}
					\item[Online Access] \href{https://www.josephsmithpapers.org/paper-summary/minutes-16-january-1836/1}{Here}
					% \item[Analysis] \hyperref[XX]{Here}
				\end{description}
				
				\begin{description}
					\item[Background]
				\end{description}
				
					% It might also be useful to give a two sentence context of the period: e.g., "This speech was given in Nauvoo to a small congregation one month prior to the destruction of the Nauvoo Expositor. These and other tensions may have been on the prophet's mind when he gave the speech."
				
				\begin{description}
					\item[Original Source]
				\end{description}
				
					% brief description of original source material, with reference to JSP or somewhere else for more info
					% Background material: it might be helpful to have a note that says whether a given document was presented as a speech, was written in a journal, etc.
					% Also a comment here about how reliable the source might be, or what shape it is in, if there are large omissions or parts that are hard to understand, and so on.
				
				\begin{description}
					\item[Text]
				\end{description}
				
						\hypertarget{EMS-1836-JS-16-JAN-1836-a}%
					Saturday
						\marginnote{a}%
					morning the 16\supers{th} by request I meet with the council of the 12 in company with my colleagues F[rederick] G Williams and S[idney] Rigdon

					Council organized and opened by singing and prayer offered up by Thomas B. Marsh president of the 12

					He arose and requested the privilege in behalf of his colleagues of speaking, each in his turn without being interupted; which was granted them--- Elder Marsh proceeded to unbosom his feelings touching the mission of the 12, and more particularly respecting a certain letter which they recieved from the presidency of the high council in Kirtland, while attending a conference in the \sout{East} State of Maine--- also spoke of being plased in our council, on Friday last below the council's of Kirtland and Zion having been previously placed next [to] the presidency, in our assemblies--- also observed that they were hurt on account of some remarks made by President H[yrum] Smith on the trial of Gladden Bishop who had been previously tried before the council of the 12, while on their mission in the east, who had by their request thrown his case before the high council in Kirtland for investigation, and the 12 concidered that their proceedings with him were in some degree, discountenanced---

						\hypertarget{EMS-1836-JS-16-JAN-1836-b}%
					\phantom{ }\sout{The}
						\marginnote{b}%
					\sout{rest} \sout{<remaining>} Elder Marsh then gave way to his brethren and they arose and spoke in turn untill they had all spoken acquiessing in the observations of Elder Marsh and mad[e] some additions to his remarks which are as follows--- That the letter in question which they received from the presidency, in which two of their numbers were suspended, and the rest severely chastened, and that too upon testimony which was unwarantable, and particular\sout{ly} stress was laid upon a certain letter which the presidency had received from Dr. W[arren] A. Cowdery of Freedom New York in which he prefered charges against them which were false, and upon which \sout{they} <we> (the presiders) had acted in chastning them and therefore, the 12, had concluded that the presidency had lost confidence in them, and that whereas the church in this place, had carressed them, at the time of their appointment, to the appostleship they now treated them coolly and appear to have lost confidence in them also---

						\hypertarget{EMS-1836-JS-16-JAN-1836-c}%
					They
						\marginnote{c}%
					spoke of their having been in this work from the beginning almost and had born the burden in the heat of the day and passed through many trials and that the presidency ought not to \sout{have} suspect their fidelity nor loose confidence in them, neither have chastised them upon such testimony as was lying \sout{before} before them--- also urged the necessity of an explanition upon the letter which they received from the presidency, and the propriety of their having information as it respects their duties, authority \&c--- that they might come to <an> understanding in all things, that they migh[t] act in perfect unison and harmony before the Lord and be prepared for the endument--- also that they had prefered a charge against Dr [Warren A.] Cowdery for his unchristian conduct which the presidency had disregarded--- also that President O[liver] Cowdery on a certain occasion had made use of language to one of the twelve that was unchristian and unbecoming any man, and that they would not submit to such treatment

						\hypertarget{EMS-1836-JS-16-JAN-1836-d}%
					The
						\marginnote{d}%
					remarks of all the 12 were made in a verry forcible and explicit manner yet cool and deliberate; I arose

					I observed that we had heard them patiently and in turn should expect to be heard patiently also; and first I remarked that it was necessary that the 12 should state whether they were determined to persevere in the work of the Lord, whether the presidency are able to satisfy them or not; vote called and carried in the affirmative unaminously; I then said to them that I had not lost confidence in them, and that they had no reason to suspect my confidence, and that I would be willing to be weighed in the scale of truth today in this matter, and risk it in the day of judgment; and as it respects the chastning contained in the letter in question which I acknowledge might have been expressed in too harsh language; which was not intentional and I ask your forgiveness in as much as I have hurt your feelings; but nevertheless, the letter \sout{that} that Elder Mc.lellen [William E. McLellin] wrote back to Kirtland while the twelve were at the east was harsh also and I was willing to set the one against the other;
						\hypertarget{EMS-1836-JS-16-JAN-1836-e}%
					I
						\marginnote{e}%
					next proceeded to explain the subject of the duty of the twelve; and their authority which is next to the present presidency, and that the arangement of the assembly in this place on the 15 \supers{inst} in placing the high councils of Kirtland \sout{and} next [to] the presidency was because the buisness to be transacted was buisness that related to that body in particular which was to fill the several quorum's in Kirtland; not beca[u]se they were first in office, and that the arangement was most Judicious that could be made on the occassion also the 12, are not subject to any other than the first presidency; viz. myself S. Rigdon and F. G. Williams--- I also stated to the 12, that I do not \sout{continue} countinance the harsh language of President [Oliver] Cowdery to them neither in myself nor any other man, although I have sometimes spoken to[o] harsh from the impulse of the moment and inasmuch as I have wounded your feelings brethren I ask your forgivness,
						\hypertarget{EMS-1836-JS-16-JAN-1836-f}%
					for
						\marginnote{f}%
					I love you and will hold you up with all my heart in all righteousness before the Lord, and before all men, for be assured brethren I am willing to stem the torrent of all opposition; in storms in tempests in thunders and lightning by sea and by land in the wilderness or among fals brethren, or mobs or wherever God in his providence may call us and I am determined that neither hights nor depths principalities nor powers things present or to come nor any other creature shall separate me from you; and I will now covenant with you before God that I will not listen too nor credit, any derogatory report against any of you nor condemn you upon any testimony beneath the heavens, short of that testimony which is infalible, untill I can see you face to face and know of a surity and I do place unlimited confidence in your word for I believe you to be men of truth, and I ask the same of you, when I tell you any thing that you place equal confidence in my word for I will not tell you I know anything which I do not know--- but I have already consumed more time than I intended to when I commenced and I will now give way to my colleagues

						\hypertarget{EMS-1836-JS-16-JAN-1836-g}%
					President
						\marginnote{g}%
					Rigdon arose next and acquiessed in what I had said and acknowledged to the 12, that he had not done as he ought, in not citing Dr. Cowdery to trial on the charges that were put into his hands by the 12, that he had neglected his duty in this thing, for which he asked their forgiveness, and would now attend to it if they desired him to do so, and \sout{Elder} <Presdt> Rigdon also observed to the 12 \sout{that} <if he> \sout{he} \sout{might} \sout{have} <had> spoken, or reproved too \sout{harshe,} <harshly,> at any time and had injured their feelings by so doing he asked their forgivness.---

					President Williams arose and acquiessed in the above sentiment's expressed by myself and President Rigdon, in full and said many good things

						\hypertarget{EMS-1836-JS-16-JAN-1836-h}%
					The
						\marginnote{h}%
					President of the 12, then called a vote of that body to know whether they were perfectly satisfied with the explenation which we had given them and whether they would enter into the covenant we had proposed to them, which was most readily manifested in the affirmative by raising their hands to heaven, in testimony of their willingness and desire to enter into this covenant and their entire satisfaction with our explanation, upon all the difficulties that were on their minds, we then took each others by the hand in confirmation of our covenant and their was a perfect unison of feeling on this occasion, and our hearts overflowed with blessings, which we pronounced upon eachothers heads as the Spirit gave us utterance my scribe is included in this covenant \sout{with} and blessings with us, for I love him, for <the> truth and integrity that dwelleth in him and may God enable us all, to perform our vows and covenants with each other in all fidelity and rightiousness before Him, that our influence may be felt among the nations of the earth in mighty power, even to rend the kingdom of darkness in sunder, and triumph over priestcraft and spiritual wickedness in high places, and brake in pieces all \sout{other} kingdoms that are opposed to the Kingdom of Christ, and spread the light and truth of the everlasting gospel from the rivers to the ends of the earth

					Elder Beemon [Alvah Beman] call[ed] for council upon the subject of his returning home he wished to know whether it was best for him to return before the Solemn Assembly or not, after taking it into concideration the council advised him to tarry we dismissed by singing and prayer and retired
					
					\begin{tightright}
					W[arren] Parrish \uline{scribe}
					\end{tightright}
					
		% -----
		\subsubsection{Record, 17 January 1836}
		% -----
			\label{EMS-1836-JS-17-JAN-1836}
			% \label{EMS-1830-JS-DAY-3LETTERMONTH-YEAR}
			
			\begin{description}
					\item[Print Sources]
				\end{description}

					\begin{tabular}{p{0.5in}p{1in}p{0.5in}p{1in}}
					JSP	 	& J1:161		& JSR 		& --- \\
					DC	 	& ---				& HC 		& --- \\
					HDDC 	& ---				& TCHC 		& --- \\
					PJS 	& 1:207			& 	 		&  \\
					\end{tabular}
					% HDDC = woodford's DC diss; JSR = Marquardt's JS Revelations; TCHC = Vogel HC
				
				\begin{description}
					\item[Online Access] \href{https://www.josephsmithpapers.org/paper-summary/journal-1835-1836/128}{Here}
					% \item[Analysis] \hyperref[XX]{Here}
				\end{description}
				
				\begin{description}
					\item[Background]
				\end{description}
				
					% It might also be useful to give a two sentence context of the period: e.g., "This speech was given in Nauvoo to a small congregation one month prior to the destruction of the Nauvoo Expositor. These and other tensions may have been on the prophet's mind when he gave the speech."
				
				\begin{description}
					\item[Original Source]
				\end{description}
				
					% brief description of original source material, with reference to JSP or somewhere else for more info
					% Background material: it might be helpful to have a note that says whether a given document was presented as a speech, was written in a journal, etc.
					% Also a comment here about how reliable the source might be, or what shape it is in, if there are large omissions or parts that are hard to understand, and so on.
				
				\begin{description}
					\item[Text]
				\end{description}
				
					...we were then invited to Elder Cahoons to a feast which was prepared on the occasion, and had a good time while partaking of the rich repast that was spread before us, and I virely [verily] realized that it was good for brethren to dwell together in unity, \sout{even} like the dew upon the mountains of Israel, where the Lord command[e]d blessings, even life for ever more...
					
		% -----
		\subsubsection{Record, 20 January 1836}
		% -----
			\label{EMS-1836-JS-20-JAN-1836}
			% \label{EMS-1830-JS-DAY-3LETTERMONTH-YEAR}
			
			\begin{description}
					\item[Print Sources]
				\end{description}

					\begin{tabular}{p{0.5in}p{1in}p{0.5in}p{1in}}
					JSP	 	& J1:165--166		& JSR 		& --- \\
					DC	 	& ---				& HC 		& --- \\
					HDDC 	& ---				& TCHC 		& --- \\
					\end{tabular}
					% HDDC = woodford's DC diss; JSR = Marquardt's JS Revelations; TCHC = Vogel HC
				
				\begin{description}
					\item[Online Access] \href{https://www.josephsmithpapers.org/paper-summary/journal-1835-1836/133}{Here}
					% \item[Analysis] \hyperref[XX]{Here}
				\end{description}
				
				\begin{description}
					\item[Background]
				\end{description}
				
					% It might also be useful to give a two sentence context of the period: e.g., "This speech was given in Nauvoo to a small congregation one month prior to the destruction of the Nauvoo Expositor. These and other tensions may have been on the prophet's mind when he gave the speech."
				
				\begin{description}
					\item[Original Source]
				\end{description}
				
					% brief description of original source material, with reference to JSP or somewhere else for more info
					% Background material: it might be helpful to have a note that says whether a given document was presented as a speech, was written in a journal, etc.
					% Also a comment here about how reliable the source might be, or what shape it is in, if there are large omissions or parts that are hard to understand, and so on.
				
				\begin{description}
					\item[Text]
				\end{description}
				
						\hypertarget{EMS-1836-JS-20-JAN-1836-a}%
					...after
						\marginnote{a}%
					the above arangments were made Eld\supers{r.} Boynton \& his Lady with their attendants, came in and were seated in front of the presidency,--- a hymn was sung, after which I adressed a throne of grace,--- I then arose and read aloud \sout{the} <a> licence granting any minister of the gospel the priviledge of solemnizing the rights of matrimony, and after calling for objection if any there were, against the anticipated alliance between Eldr. Boynton \& Miss Lowell \sout{to} and waiting sufficient time, I observed that all forever after this must hold their peace---
					
						\hypertarget{EMS-1836-JS-20-JAN-1836-b}%
					I
						\marginnote{b}%
					then envited them to join hands and I pronounced the ceremony according to the rules and regulations of the church of the Latter-day-Saints, \sout{and} <in> the name of God, and in the name of Jesus Christ, I pronounced upon them the blessings of Abraham Isaac and Jacob and such other blessings as the Lord put into my heart, and being much under the influence of a cold I then gave way and President S[idney] Rigdon arose and delivered a verry forcible address, suited to the occasion, and closed the servises of the evening by prayer--- Eldr. O. Hyde Eld\supers{r.} L[uke] Johnson \& Eldr. W. Parrish who served on the occasion, then presented the presidency with three servers filled with glasses of wine, to bless, and it fell to my lot to attend to this duty, which I cheerfully discharged, it was then passed round in order, then the cake, in the same order, and suffise it to say our hearts were made \sout{cheerful} \sout{and} glad, while partaking of the bounty of the earth, which was presented, untill we had taken our fill, and Joy filled every bosom, and the countenances of old, and young, alike, seemed to bloom with the cheerfulness and smiles of youth and an entire unison of feeling seemed to pervade the congregation, and indeed I doubt whether the pages of history can boast of a more splendid and inocent wedding and feast than this for it was conducted after the order of heaven, who has a time for all thing and this being a time of rejoicing, we hartily embraced it, and conducted ourselves accordingly...
					
		% -----
		\subsubsection{Visions and Records, 21 January 1836 [D\&C 137]}
		% -----
			\label{EMS-1836-JS-21-JAN-1836}
			% \label{EMS-1830-JS-DAY-3LETTERMONTH-YEAR}
			
			\begin{description}
					\item[Print Sources]
				\end{description}

					\begin{tabular}{p{0.5in}p{1in}p{0.5in}p{1in}}
					JSP	 	& D5:157--160		& JSR 		& 278--279 \\
					DC	 	& Section 137		& HC 		& 2:380--382 \\
					HDDC 	& ---				& TCHC 		& 2:372--373 \\
					TPJS 	& 106--108			& 	 		&  \\
					\end{tabular}
					% HDDC = woodford's DC diss; JSR = Marquardt's JS Revelations; TCHC = Vogel HC
				
				\begin{description}
					\item[Online Access] \href{https://www.josephsmithpapers.org/paper-summary/visions-21-january-1836-dc-137/1}{Here}
					% \item[Analysis] \hyperref[DC137]{Here}
				\end{description}
				
				\begin{description}
					\item[Background]
				\end{description}
				
					% It might also be useful to give a two sentence context of the period: e.g., "This speech was given in Nauvoo to a small congregation one month prior to the destruction of the Nauvoo Expositor. These and other tensions may have been on the prophet's mind when he gave the speech."
				
				\begin{description}
					\item[Original Source]
				\end{description}
				
					% brief description of original source material, with reference to JSP or somewhere else for more info
					% Background material: it might be helpful to have a note that says whether a given document was presented as a speech, was written in a journal, etc.
					% Also a comment here about how reliable the source might be, or what shape it is in, if there are large omissions or parts that are hard to understand, and so on.
				
				\begin{description}
					\item[Text]
				\end{description}
				
						\hypertarget{EMS-1836-JS-21-JAN-1836-a}%
					Thursday
						\marginnote{a}%
					morning the 21\supers{st} This morning a minister from conneticut by the name of John W. Olived called at my house and enquired of my father if \sout{Smith} the pro[p]het live's here he replied that he did not understand him. Mr. Olived asked the same question again and again and recieved the same answer, he finally asked if Mr. Smith lives here, father replyed O yes Sir I understand you now,--- father then stept into my room, and informed me that a gentleman had called to see me, I went into the room where he was,
						\hypertarget{EMS-1836-JS-21-JAN-1836-b}%
					and
						\marginnote{b}%
					the first question he asked me, after passing a compliment, was to know how many members we have in our church, I replyed to him, that we hav \sout{about} between 15 hundred and 2,000 in this branch,--- He then asked me wherin we differ from other christian denomination I replyed that we believe the bible, and they do not,--- however he affirmed that he believed the bible, I told him then to be baptised,---
						\hypertarget{EMS-1836-JS-21-JAN-1836-c}%
					he
						\marginnote{c}%
					replied that he did not realize it to be his duty--- But \sout{after} <when> \sout{laying} <laid> \sout{him} before him the principles of the gospel, viz. faith and repentance and baptism for the remission <of sins> and the laying on of hands for the reseption of the Holy Ghost <he manifested much surprise>--- I then observed that the hour for school had arived, and I must attend The man seemed astonished at our doctrine but by no means hostile
					
						\hypertarget{EMS-1836-JS-21-JAN-1836-d}%
					At
						\marginnote{d}%
					about 3, oclock P.M I dismissed the School and the presidency; retired to the loft of the printing office, where we attended to the ordinance of washing our bodies in pure water, we also perfumed our bodies and our heads, in the name of the Lord at early candlelight, I meet with the presidency, at the west school room in the Chapel to attend to the ordinance of annointing our heads with holy oil---
						\hypertarget{EMS-1836-JS-21-JAN-1836-e}%
					also
						\marginnote{e}%
					the councils of \sout{Zion} Kirtland and Zion, meet in the two adjoining rooms, who waited in prayer while we attended to the ordinance,--- I took the oil in my <left> \sout{right} hand, father Smith being seated before me and the rest of the presidency encircled him round about,--- we then streched our right hands to heaven and blessed the oil and concecrated it in the name of Jesus Christ---
						\hypertarget{EMS-1836-JS-21-JAN-1836-f}%
					we
						\marginnote{f}%
					then laid our hands on, our aged fath[er] Smith, and invoked, the blessings of heaven,--- I then annointed his head with the concecrated oil, and sealed many blessings upon \sout{his} <him,> \sout{head,} the presidency then in turn, laid their hands upon his head, begenning at the eldest, untill they had all laid their hands on him, and pronounced such blessings, upon his head as the Lord put into their hearts--- all blessing him to be our patraark [patriarch], \sout{and} <to> annoint our heads, and attend to all duties that pertain to \sout{this} <that> office.---
						\hypertarget{EMS-1836-JS-21-JAN-1836-g}%
					I
						\marginnote{g}%
					then took the seat, and father annoint[ed] my head, and sealed upon me, the blessings, of Moses, to lead Israel in the latter days, even as moses led \sout{them} <him> in days of old,--- also the blessings of Abraham Isaac and Jacob,--- all of the presidency laid their hands upon me and pronounced upon my head many prophesies, and blessings, many of which I shall not notice at this time, but as Paul said, so say I, let us come to vissions and revelations, \sout{the}---
						\hypertarget{EMS-1836-JS-21-JAN-1836-h}%
					The
						\marginnote{h}%
					heavens were opened upon us and I beheld the celestial kingdom of God, and the glory thereof, whether in the body or out I cannot tell,--- I saw the transcendant beauty of the gate \sout{that} \sout{enters}, through which the heirs of that kingdom will enter, which was like unto circling flames of fire, also the blasing throne of God, whereon was seated the Father and the Son,--- I saw the beautiful streets of that kingdom, which had the appearance of being paved with gold---
						\hypertarget{EMS-1836-JS-21-JAN-1836-i}%
					I
						\marginnote{i}%
					saw father Adam, and Abraham and Michael and my father and mother, my brother Alvin that has long since slept, and marvled how it was that he had obtained \sout{this} an inheritance <in> \sout{this} <that> kingdom, seeing that he had departed this life, before the Lord <had> set his hand to gather Israel <the second time> and had not been baptized for the remission of sins---
						\hypertarget{EMS-1836-JS-21-JAN-1836-j}%
					Thus
						\marginnote{j}%
					\phantom{ }\sout{said} came the voice <of the Lord un>to me saying all who have died with[out] a knowledge of this gospel, who would have received it, if they had been permited to tarry, shall be heirs of the celestial kingdom of God--- also all that shall die henseforth, with<out> a knowledge of it, who would have received it, with all their hearts, shall be heirs of that kingdom, for I the Lord <will> judge all men according to their works according to the desires of their hearts---
						\hypertarget{EMS-1836-JS-21-JAN-1836-k}%
					and
						\marginnote{k}%
					\phantom{ }\sout{again} \sout{I} \sout{also} \sout{beheld} \sout{the} \sout{Terrestial} \sout{kingdom} I also beheld that all children who die before they arive to the years of accountability, are saved in the celestial kingdom of heaven--- I saw the 12, apostles of the Lamb, who are now upon the earth who hold the keys of this last ministry, in foreign lands, standing together in a circle much fatiegued, with their clothes tattered and feet swolen, with their eyes cast downward, and Jesus <standing> in their midst, and they did not behold him, \sout{he} the Saviour looked upon them and wept---
						\hypertarget{EMS-1836-JS-21-JAN-1836-l}%
					I
						\marginnote{l}%
					also beheld Elder McLellen [William E. McLellin] in the south, standing upon a hill surrounded with a vast multitude, preaching to them, and a lame man standing before him, supported by his crutches, he threw them down at his word, and leaped as an heart [hart] by the mighty power of God
					
						\hypertarget{EMS-1836-JS-21-JAN-1836-m}%
					Also
						\marginnote{m}%
					Eldr Brigham Young standing in a strange land, in the far southwest, in a desert place, upon a rock in the midst of about a dozen men of colour, who, appeared hostile He was preaching to them in their own toung, and the angel of God standing above his head with a drawn sword in his hand protecting him, but he did not see it,---
						\hypertarget{EMS-1836-JS-21-JAN-1836-n}%
					and
						\marginnote{n}%
					I finally saw the 12 in the celestial kingdom of God,--- I also beheld the redemption of Zion, and many things which the toung of man, cannot discribe in full, ,--- Many of my brethren who received this ordinance with me, saw glorious visions also,--- angels ministered unto them, as well as my self, and the power of the highest rested upon, us the house was filled with the glory of God, and we shouted Hosanah to \sout{the} God and the Lamb
					
						\hypertarget{EMS-1836-JS-21-JAN-1836-o}%
					I
						\marginnote{o}%
					am mistaken, concerning my receiving the holy anointing first after father Smith, we received <it> in turn according to our age, (that is the presidency,)
					
					My Scribe also recieved his anointing <with us> and saw in a vision the armies of heaven protecting the Saints in their return to Zion--- <\& many things that I saw>
					
						\hypertarget{EMS-1836-JS-21-JAN-1836-p}%
					The
						\marginnote{p}%
					Bishop of Kirtland with his counsellors and the Bishop of Zion with his counsellors, were present with us, and received their, annointing under the hands of father Smith and confirmed by the presidency and the glories of heaven was unfolded to them also---
					
					We then invited the counsellors of Kirtland and Zion \sout{and} \sout{Kirtland} into our room, and President Hyrum Smith annointed the head of the president of the counsellors in Kirtland and President D[avid] Whitmer the head of the president, of the counsellors of Zion---
					
					The president of each quorum then annointed the heads of his colleagues, each in his turn beginning, at the eldest
					
						\hypertarget{EMS-1836-JS-21-JAN-1836-q}%
					The
						\marginnote{q}%
					vision of heaven \sout{were} <was> opened to these also, some of them saw the face of the Saviour, and others were ministered unto by holy angels, and the spirit of propesey and revelation was poured out in mighty power, and loud hosanahs and glory to God in the highest, saluted the heavens for we all communed with the h[e]avenly host's,--- and I saw in my vision all of the presidency in the Celistial Kingdom of God, and, many others who were present
					
					Our meeting was opened by singing and prayer offered up by the head of each quorum, and closed by singing and invoking the benediction of heaven with uplifted hands, and retired between one and 2, oclock in the morning [\textit{9 lines blank}]
					
		% -----
		\subsubsection{Record, 22 January 1836}
		% -----
			\label{EMS-1836-JS-22-JAN-1836}
			% \label{EMS-1830-JS-DAY-3LETTERMONTH-YEAR}
			
			\begin{description}
					\item[Print Sources]
				\end{description}

					\begin{tabular}{p{0.5in}p{1in}p{0.5in}p{1in}}
					JSP	 	& ---		& JSR 		& --- \\
					DC	 	& ---		& HC 		& --- \\
					HDDC 	& ---		& TCHC 		& --- \\
					TPJS 	& ---		& 	 		&  \\
					\end{tabular}
					% HDDC = woodford's DC diss; JSR = Marquardt's JS Revelations; TCHC = Vogel HC
				
				\begin{description}
					\item[Online Access] \href{https://www.josephsmithpapers.org/paper-summary/journal-1835-1836/141}{Here}
					% \item[Analysis] \hyperref[DC137]{Here}
				\end{description}
				
				\begin{description}
					\item[Background]
				\end{description}
				
					% It might also be useful to give a two sentence context of the period: e.g., "This speech was given in Nauvoo to a small congregation one month prior to the destruction of the Nauvoo Expositor. These and other tensions may have been on the prophet's mind when he gave the speech."
				
				\begin{description}
					\item[Original Source]
				\end{description}
				
					% brief description of original source material, with reference to JSP or somewhere else for more info
					% Background material: it might be helpful to have a note that says whether a given document was presented as a speech, was written in a journal, etc.
					% Also a comment here about how reliable the source might be, or what shape it is in, if there are large omissions or parts that are hard to understand, and so on.
				
				\begin{description}
					\item[Text]
				\end{description}
				
						\hypertarget{EMS-1836-JS-22-JAN-1836-a}%
					Friday
						\marginnote{a}%
					morning the 22\supers{ond} attended at the school room at the us[u]al hour,--- But insted of persuing our studies we \sout{commenced} spent the time in rehearsing to each other the glorious scenes that transpired on the preceding evening, while attending to the ordinance of holy anointing.--- At evening we met at the same place, with the council of the 12 and the presidency of the 70 who were to receive this ordinance; the high councils of Kirtland and Zion were present also: we called to order and organized; the Presidency then proceeded to consecrate the oil; we then laid our hands upon Elder Thomas B. Marsh who is the president of the 12 and ordained him to the authority of anointing his brethren, I then poured the concecrated oil upon his head in the name of Jesus Christ and sealed such blessings upon him as the Lord put into my heart; the rest of the presidency then laid their hands upon him and blessed him each in their turn beginning at the eldest; he then anointed <and blessed> his brethren from the oldest to the youngest, I also laid my hands upon them and prounounced many great and glorious upon their heads; the heavens were opened and angels ministered unto us.
					
						\hypertarget{EMS-1836-JS-22-JAN-1836-b}%
					The\sout{n}
						\marginnote{b}%
					12 then proceeded to anoint and bless the presidency of the 70 and seal upon their heads power and authority to anoint their brethren; the heavens were opened upon Elder Sylvester Smith and he leaping up exclaimed, The horsemen of Israel and the chariots thereof. \sout{President} \sout{Rigdon} \sout{arose} Br. Carloss [Don Carlos] Smith was also, annointed and \sout{ordained} blessed to preside over the high preisthood.--- President [Sidney] Rigdon, arose to conclude the servises of the evening by invoking the benediction of heaven \sout{of heaven} upon the Lords anointed <which he did> in an eloquent manner the congregation shouted a loud hosannah the gift of toungs, fell upon us in mighty power, angels mingled \sout{themselves} their voices with ours, while their presence was in our midst, and unseasing prasis [praises] swelled our bosoms for the space of half an hour,--- I then observed to the brethren that it was time to retire, we accordingly <closed> our interview and returned home at about 2. oclock in the morning \& the spirit \& visions of God attended me through the night
					
		% -----
		\subsubsection{Record, 24 January 1836}
		% -----
			\label{EMS-1836-JS-24-JAN-1836}
			% \label{EMS-1830-JS-DAY-3LETTERMONTH-YEAR}
			
			\begin{description}
					\item[Print Sources]
				\end{description}

					\begin{tabular}{p{0.5in}p{1in}p{0.5in}p{1in}}
					JSP	 	& ---		& JSR 		& --- \\
					DC	 	& ---		& HC 		& --- \\
					HDDC 	& ---		& TCHC 		& --- \\
					TPJS 	& ---		& 	 		&  \\
					\end{tabular}
					% HDDC = woodford's DC diss; JSR = Marquardt's JS Revelations; TCHC = Vogel HC
				
				\begin{description}
					\item[Online Access] \href{https://www.josephsmithpapers.org/paper-summary/journal-1835-1836/142}{Here}
					% \item[Analysis] \hyperref[DC137]{Here}
				\end{description}
				
				\begin{description}
					\item[Background]
				\end{description}
				
					% It might also be useful to give a two sentence context of the period: e.g., "This speech was given in Nauvoo to a small congregation one month prior to the destruction of the Nauvoo Expositor. These and other tensions may have been on the prophet's mind when he gave the speech."
				
				\begin{description}
					\item[Original Source]
				\end{description}
				
					% brief description of original source material, with reference to JSP or somewhere else for more info
					% Background material: it might be helpful to have a note that says whether a given document was presented as a speech, was written in a journal, etc.
					% Also a comment here about how reliable the source might be, or what shape it is in, if there are large omissions or parts that are hard to understand, and so on.
				
				\begin{description}
					\item[Text]
				\end{description}
				
					In the evening met the Presidency in \sout{high} <the> room over the printing room \& counseled on the subject of endowment \& the preperation necessary for the solemn Assembly which is to be called when the House of the Lord is finished
					
		% -----
		\subsubsection{Record, 28 January 1836}
		% -----
			\label{EMS-1836-JS-28-JAN-1836}
			% \label{EMS-1830-JS-DAY-3LETTERMONTH-YEAR}
			
			\begin{description}
					\item[Print Sources]
				\end{description}

					\begin{tabular}{p{0.5in}p{1in}p{0.5in}p{1in}}
					JSP	 	& ---		& JSR 		& --- \\
					DC	 	& ---		& HC 		& --- \\
					HDDC 	& ---		& TCHC 		& --- \\
					TPJS 	& ---		& 	 		&  \\
					\end{tabular}
					% HDDC = woodford's DC diss; JSR = Marquardt's JS Revelations; TCHC = Vogel HC
				
				\begin{description}
					\item[Online Access] \href{https://www.josephsmithpapers.org/paper-summary/journal-1835-1836/144}{Here}
					% \item[Analysis] \hyperref[XX]{Here}
				\end{description}
				
				\begin{description}
					\item[Background]
				\end{description}
				
					% It might also be useful to give a two sentence context of the period: e.g., "This speech was given in Nauvoo to a small congregation one month prior to the destruction of the Nauvoo Expositor. These and other tensions may have been on the prophet's mind when he gave the speech."
				
				\begin{description}
					\item[Original Source]
				\end{description}
				
					% brief description of original source material, with reference to JSP or somewhere else for more info
					% Background material: it might be helpful to have a note that says whether a given document was presented as a speech, was written in a journal, etc.
					% Also a comment here about how reliable the source might be, or what shape it is in, if there are large omissions or parts that are hard to understand, and so on.
				
				\begin{description}
					\item[Text]
				\end{description}
				
						\hypertarget{EMS-1836-JS-28-JAN-1836-a}%
					Thursday
						\marginnote{a}%
					<28> attended school at the usual hours In the evening met the quorems of High Priests in the west room of the upper loft of the \sout{Lord,} Lord<'s> house \& in company with my council of the presidency--- consecrated \& anointed the cousellors of the President of the High priesthood \& having instructed them \& set the quorem in order I left them to perform the holy anointing--- \& went to the quorem of Elders in the other end of the room. I assisted in anointing the counsellors of the President of the Elders \& gave them the instruction necessary for the occasion \& left the President \& his council to anoint the Elders while I should go to the adjoining room \& attend to organizing \& instructing of the quorem of the Seventy---

						\hypertarget{EMS-1836-JS-28-JAN-1836-b}%
					I
						\marginnote{b}%
					found the Twelve Apostles assembled with this quorem \& I proceeded with the quorem of the presedincy to instruct them \& also the seven presidents of the seventy Elders to call upon God with uplifted hands to seal the blessings which had been promised to them by the holy anoint[in]g As I organized this quorem with the presedincy in this room, Pres. Sylvester Smith saw a piller of fire rest down \& abide upon the heads of the quorem as we stood in the midst of the Twelve.

					When the Twelve \& the seven were through with their sealing prayers I called upon Pres. S[idney] Rigdon to seal them with uplifted hands \& when he he had done this \& cried hossannah that all congregation should join him \& shout hosannah to God \& the Lamb \& glory to God in the highest--- It was done so \& Eld. Roger Orton saw a \sout{flaming} <mighty> Angel riding upon a horse of fire with a flaming sword in his hand followed by five others--- encircle the house \& protect the saints even the Lords anointed from the power of Satan \& a host of evil spirits which were striving to disturb the saints---

						\hypertarget{EMS-1836-JS-28-JAN-1836-c}%
					Pres.
						\marginnote{c}%
					W\supersu{m} Smith one of the Twelve saw the h[e]avens op[e]ned \& the Lords host protecting the Lords anointed, Pres. Z. Coltrine [Zebedee Coltrin] one of the seven saw the saviour extended before him as upon the cross \& little after crowned with a glory upon his head above the brightness of the sun after these things were over \& \sout{I} a glorious vision which I saw had passed I instructed the seven presidents to proceede \& anoint the seventy \& returned to the room of the High Priests \& Elders \& attended to the sealing of what they had done with uplifted hands, the Lord had assisted my bro. Carloss [Don Carlos Smith] the Pres. of the High Priests to go forward with the anointing of the High priests so that he had performed it to the acceptance of the Lord, notwithstanding he was verry young \& inexperienced in such duties \& I felt to praise God with a loud hossannah for his goodness to me \& my Fathers family \& to all the children of men--- praise the Lord all ye his saints--- praise his holy name--- after these quorems were dismissed I retired to my home filled with the spirit \& my soul cried hossannah to God \& the Lamb through <the> silent watches of the night \& while my eyes were closed in sleep the visions of the Lord were sweet unto me \& his glory was round about me

					praise the Lord---
					
		% -----
		\subsubsection{Record, 30 January 1836}
		% -----
			\label{EMS-1836-JS-30-JAN-1836}
			% \label{EMS-1830-JS-DAY-3LETTERMONTH-YEAR}
			
			\begin{description}
					\item[Print Sources]
				\end{description}

					\begin{tabular}{p{0.5in}p{1in}p{0.5in}p{1in}}
					JSP	 	& ---		& JSR 		& --- \\
					DC	 	& ---		& HC 		& --- \\
					HDDC 	& ---		& TCHC 		& --- \\
					TPJS 	& ---		& 	 		&  \\
					\end{tabular}
					% HDDC = woodford's DC diss; JSR = Marquardt's JS Revelations; TCHC = Vogel HC
				
				\begin{description}
					\item[Online Access] \href{https://www.josephsmithpapers.org/paper-summary/journal-1835-1836/149}{Here}
					% \item[Analysis] \hyperref[XX]{Here}
				\end{description}
				
				\begin{description}
					\item[Background]
				\end{description}
				
					% It might also be useful to give a two sentence context of the period: e.g., "This speech was given in Nauvoo to a small congregation one month prior to the destruction of the Nauvoo Expositor. These and other tensions may have been on the prophet's mind when he gave the speech."
				
				\begin{description}
					\item[Original Source]
				\end{description}
				
					% brief description of original source material, with reference to JSP or somewhere else for more info
					% Background material: it might be helpful to have a note that says whether a given document was presented as a speech, was written in a journal, etc.
					% Also a comment here about how reliable the source might be, or what shape it is in, if there are large omissions or parts that are hard to understand, and so on.
				
				\begin{description}
					\item[Text]
				\end{description}
				
					in the evening went to the upper rooms of the Lord's house \& set the different quorems in order--- instructed the Presidents of the seventy concerning the order of their anointing \& requested them to proceed \& anoint the seventy having set all the quorems in order I re turned to my house being weary with continual anxiety \& labour in puting all the Authorities in [order?] \& in striving to purify them for the solemn assembly according to the commandment of the Lord
					
		% -----
		\subsubsection{Record, 1 February 1836}
		% -----
			\label{EMS-1836-JS-1-FEB-1836}
			% \label{EMS-1830-JS-DAY-3LETTERMONTH-YEAR}
			
			\begin{description}
					\item[Print Sources]
				\end{description}

					\begin{tabular}{p{0.5in}p{1in}p{0.5in}p{1in}}
					JSP	 	& ---		& JSR 		& --- \\
					DC	 	& ---		& HC 		& --- \\
					HDDC 	& ---		& TCHC 		& --- \\
					TPJS 	& ---		& 	 		&  \\
					\end{tabular}
					% HDDC = woodford's DC diss; JSR = Marquardt's JS Revelations; TCHC = Vogel HC
				
				\begin{description}
					\item[Online Access] \href{https://www.josephsmithpapers.org/paper-summary/journal-1835-1836/150}{Here}
					% \item[Analysis] \hyperref[XX]{Here}
				\end{description}
				
				\begin{description}
					\item[Background]
				\end{description}
				
					% It might also be useful to give a two sentence context of the period: e.g., "This speech was given in Nauvoo to a small congregation one month prior to the destruction of the Nauvoo Expositor. These and other tensions may have been on the prophet's mind when he gave the speech."
				
				\begin{description}
					\item[Original Source]
				\end{description}
				
					% brief description of original source material, with reference to JSP or somewhere else for more info
					% Background material: it might be helpful to have a note that says whether a given document was presented as a speech, was written in a journal, etc.
					% Also a comment here about how reliable the source might be, or what shape it is in, if there are large omissions or parts that are hard to understand, and so on.
				
				\begin{description}
					\item[Text]
				\end{description}
				
					in the evening attended to <the> organizing \sout{the} of the quorems of High priests--- Elders--- Seventy \& Bishops in the uper rooms of the house of the Lord \& after blessing each quorem in the name of the Lord I left them \& returned home
					
					had an other interview with M\supersu{r} Seixas our hebrew teacher \& related to him some of the dealings of God to me--- \& gave him some of the evidences of the work of the latter days--- he list[e]ned candidly \& did not oppose
					
		% -----
		\subsubsection{Record, 6 February 1836}
		% -----
			\label{EMS-1836-JS-6-FEB-1836}
			% \label{EMS-1830-JS-DAY-3LETTERMONTH-YEAR}
			
			\begin{description}
					\item[Print Sources]
				\end{description}

					\begin{tabular}{p{0.5in}p{1in}p{0.5in}p{1in}}
					JSP	 	& ---		& JSR 		& --- \\
					DC	 	& ---		& HC 		& --- \\
					HDDC 	& ---		& TCHC 		& --- \\
					TPJS 	& ---		& 	 		&  \\
					\end{tabular}
					% HDDC = woodford's DC diss; JSR = Marquardt's JS Revelations; TCHC = Vogel HC
				
				\begin{description}
					\item[Online Access] \href{https://www.josephsmithpapers.org/paper-summary/journal-1835-1836/152}{Here}
					% \item[Analysis] \hyperref[XX]{Here}
				\end{description}
				
				\begin{description}
					\item[Background]
				\end{description}
				
					% It might also be useful to give a two sentence context of the period: e.g., "This speech was given in Nauvoo to a small congregation one month prior to the destruction of the Nauvoo Expositor. These and other tensions may have been on the prophet's mind when he gave the speech."
				
				\begin{description}
					\item[Original Source]
				\end{description}
				
					% brief description of original source material, with reference to JSP or somewhere else for more info
					% Background material: it might be helpful to have a note that says whether a given document was presented as a speech, was written in a journal, etc.
					% Also a comment here about how reliable the source might be, or what shape it is in, if there are large omissions or parts that are hard to understand, and so on.
				
				\begin{description}
					\item[Text]
				\end{description}
				
						\hypertarget{EMS-1836-JS-6-FEB-1836-a}%
					Saturday
						\marginnote{a}%
					6. called the anointed together to receive the seal of all their blessings, The High Priests \& Elders in the council room as usual--- The Seventy with the Twelve in the second room \& the Bishop in the 3--- I laboured with each of these quorems for some time to bring to the order which God had shown to me which is as follows--- first part to be spent in solemn prayer before god without any talking or confusion \& the conclusion with a sealing prayer by Pres. Sidney Rigdon when all the quorems are to shout with one accord a solemn hosannah to God \& the Lamb with an Amen--- amen \& amen--- \& then all take seats \& lift up their hearts in silent prayer to God \& if any obtain a prophecy or vision \sout{not} to rise \& speak that all may be edefied \& rejoice together I had considerable trouble to get all the quorems united in this order--- I went from room to room repeatedly \& charged each separately--- assuring them that it was according to the mind of God yet notwithstanding all my labour--- while I was in the east room with the Bishops quorems I f[e]lt by the spirit that something was wrong in the quorem of Elders in the west room--- \& I immediately requested Pres. O[liver] Cowdery \& H[yrum] Smith to go in \& see what was the matter---
						\hypertarget{EMS-1836-JS-6-FEB-1836-b}%
					The
						\marginnote{b}%
					quorem of Elders had not observed the order which I had given them \& were reminded of it by Pres. Carloss [Don Carlos] Smith \& mildly requested to observe order \& continue in prayer \& requested--- some of them replied that they had a teacher of their own \& did not wish to be troubled by others this caused the spirit of the Lord to withdraw This interrupted the meeting \& this quorem lost th[e]ir blessing in a great measurs--- the other quorems were more careful \& the quorem of the Seventy enjoyed a great flow of the holy spirit many arose \& spok[e] testifying that they were filled with the holy spirit which was like fire in their bones so that they could not hold their peace but were constrained to cry hosannah to God \& the Lamb \& glory in the highest. Pres. W\supersu{m} Smith one of the twelve saw a vision of the Twelve \& Seven in council together in old England \& prophecied that a great work would be done by them in the old \sout{contry} contries \& God was already beginning to work in the hearts of the p[e]ople--- Prs. Z. Coltrine [Zebedee Coltrin] one of the seven saw a vision of the Lords Host--- \& others were filled with the spirit \& spake in tongues \& prophecied--- This was a time of rejoicing long to be rememberd! praise the Lord---
					
		% -----
		\subsubsection{Blessing to Alvin Winegar, 7 February 1836}
		% -----
			\label{EMS-1836-JS-7-FEB-1836}
			% \label{EMS-1830-JS-DAY-3LETTERMONTH-YEAR}
			
			\begin{description}
					\item[Print Sources]
				\end{description}

					\begin{tabular}{p{0.5in}p{1in}p{0.5in}p{1in}}
					JSP	 	& D5:165--168		& JSR 		& --- \\
					DC	 	& ---				& HC 		& 2:393 \\
					HDDC 	& ---				& TCHC 		& 2:383 \\
					\end{tabular}
					% HDDC = woodford's DC diss; JSR = Marquardt's JS Revelations; TCHC = Vogel HC
				
				\begin{description}
					\item[Online Access] \href{https://www.josephsmithpapers.org/paper-summary/blessing-to-alvin-winegar-7-february-1836/1}{Here}
					% \item[Analysis] \hyperref[XX]{Here}
				\end{description}
				
				\begin{description}
					\item[Background]
				\end{description}
				
					% It might also be useful to give a two sentence context of the period: e.g., "This speech was given in Nauvoo to a small congregation one month prior to the destruction of the Nauvoo Expositor. These and other tensions may have been on the prophet's mind when he gave the speech."
				
				\begin{description}
					\item[Original Source]
				\end{description}
				
					% brief description of original source material, with reference to JSP or somewhere else for more info
					% Background material: it might be helpful to have a note that says whether a given document was presented as a speech, was written in a journal, etc.
					% Also a comment here about how reliable the source might be, or what shape it is in, if there are large omissions or parts that are hard to understand, and so on.
				
				\begin{description}
					\item[Text]
				\end{description}
				
						\hypertarget{EMS-1836-JS-7-FEB-1836-a}%
					A
						\marginnote{a}%
					blessing--- Pronounced by the first Presidency upon the head of Alvin Winegar--- F\supersu{eby} 7\supersu{th} 1836.
					
					We lay our hands upon thy head in the name of Jesus and ask our Heavenly Father to bestow many blessings upon thee, both in this life--- and in that which is to come--- we seal many blessings in time upon thee, and in Eternity an exceeding rich reward. Thou wast willing to lay down thy life for thy brethren and the Lord shall cho[o]se thee for his own--- Even as a vessel in the House of God--- The Lord shall seal blesings for thee in Heaven--- and look upon them and none shall fail--- thy mind shall be strengthened--- faith encreased--- receive much wisdom--- power and understanding--- must offer all that thou hast a sacrifice to thy God, even all thy mind and strengt[h]--- and give thy self up to him without reserve. body and spirit--- go at his command and he will loose thy tongue and make thee a swift Messenger--- an instrument of much good on the Earth--- see great disolations come upon the wicked and glory to the righteous--- and his power manifested in blessings and in cureings--- Shall be lifted up and shout hosannah among the heavently hosts--- see the Heavens opened.
						\hypertarget{EMS-1836-JS-7-FEB-1836-b}%
					The
						\marginnote{b}%
					charriots of Israel and the Horsemen thereof--- and great shall be thy rejoicings--- Shall receive the holy Priesthood in due time, which shall never be taken from thee in time or in Eternity--- Stand upon the Earth till the son of man comes--- see the end of this generation--- do much good in proclaiming the gospel carried to Nations afar off--- preach in all languages and tongues--- power to resist the power of floods and flames--- overcome the destroyer--- and the pestilence shall not harm thee--- power to heal the sick--- open the blind eyes--- unstop the deaf ears--- Cause the \sout{dumb} Tongue of the dumb to sing for joy--- many shall seek to touch thy garments--- and others shall send handkerchiefs, and Aprons to thee and be healed by this means. Thou shalt have all the power that thou needest to fill thy Ministry[.]
						\hypertarget{EMS-1836-JS-7-FEB-1836-c}%
					Angels
						\marginnote{c}%
					shall minister unto thee--- and thy saviour shall stand before thee, and thou shalt be able to Certify to the Nations of things which thou hast seen, and doest most assuredly know. The Heavenly Hosts shall be round about thee, to deliver thee out of all thy afflictions--- and thy soul shall Shout Hosannah to God and the Lamb and nothing shall prevail against thee to harm thee--- These blessings dear brother if thou art faithful Shall rest upon thy head and be <all> fulfilled--- for we seal them upon the[e] in the name of Jesus by the authority of the Holy Priesthood even so Amen.------ Kertland Ohio------ Sylvester Smith / scribe. JW\supers{sr} last scribe.
					
		% -----
		\subsubsection{Minutes, 12 February 1836}
		% -----
			\label{EMS-1836-JS-12-FEB-1836}
			% \label{EMS-1830-JS-DAY-3LETTERMONTH-YEAR}
			
			\begin{description}
					\item[Print Sources]
				\end{description}

					\begin{tabular}{p{0.5in}p{1in}p{0.5in}p{1in}}
					JSP	 	& D5:170--173		& JSR 		& --- \\
					DC	 	& ---				& HC 		& 2:394--395 \\
					HDDC 	& ---				& TCHC 		& 2:384--385 \\
					TPJS 	& 108				& 	 		&  \\
					\end{tabular}
					% HDDC = woodford's DC diss; JSR = Marquardt's JS Revelations; TCHC = Vogel HC
				
				\begin{description}
					\item[Online Access] \href{https://www.josephsmithpapers.org/paper-summary/minutes-12-february-1836/1}{Here}
					% \item[Analysis] \hyperref[XX]{Here}
				\end{description}
				
				\begin{description}
					\item[Background]
				\end{description}
				
					% It might also be useful to give a two sentence context of the period: e.g., "This speech was given in Nauvoo to a small congregation one month prior to the destruction of the Nauvoo Expositor. These and other tensions may have been on the prophet's mind when he gave the speech."
				
				\begin{description}
					\item[Original Source]
				\end{description}
				
					% brief description of original source material, with reference to JSP or somewhere else for more info
					% Background material: it might be helpful to have a note that says whether a given document was presented as a speech, was written in a journal, etc.
					% Also a comment here about how reliable the source might be, or what shape it is in, if there are large omissions or parts that are hard to understand, and so on.
				
				\begin{description}
					\item[Text]
				\end{description}
				
						\hypertarget{MS-1836-JS-12-FEB-1836-a}%
					Friday
						\marginnote{a}%
					evening, February 12\supersu{th}\supers{.} 1836

					Council convened in the House of the Lord for the purpose of taking into consideration the subject of ordaining men to the office of Elder and other offices in the Church of Christ of Latter Day Saints.

					Opened by singing and prayer

					The following resolutions were offered by the Presidency for discussion.

					1\supersu{st} Resolved that no one be ordained to any office in the Church in this stake of Zion at Kirtland without the unanimous voice of the several bodies that constitute this quorum, who are appointed to do church business in the name of \sout{the} <said> church. (Viz) The Presidency of the Church and Council. The twelve Apostles of the Lamb. The twelve High Counsellors of Kirtland. The twelve High Counsellors--- of Zion, The Bishop of Kirtland, and his counsellors The Bishop of Zion and his counsellors, and the seven Presidents of Seventies, until otherwise ordered by said quorum.

						\hypertarget{MS-1836-JS-12-FEB-1836-b}%
					2\supersu{dly.}
						\marginnote{b}%
					Resolved that none be ordained in the branches of said church abroad, unless they are recommended by the voice of the respective branches of the church to which they belong to the general conference appointed by the heads of the Church, and from that conference receive their ordination.

					The foregoing were concurred in by the Presidents of the Seventies.

					\begin{tightright}
					W[arren] Parrish Clerk.
					\end{tightright}
					
		% -----
		\subsubsection{Discourse, 12 February 1836}
		% -----
			\label{EMS-1836-JS-12-FEB-1836-DIS}
			% \label{EMS-1830-JS-DAY-3LETTERMONTH-YEAR}
			
			\begin{description}
					\item[Print Sources]
				\end{description}

					\begin{tabular}{p{0.5in}p{1in}p{0.5in}p{1in}}
					JSP	 	& J1:184--185		& JSR 		& --- \\
					DC	 	& ---				& HC 		& 2:394 \\
					HDDC 	& ---				& TCHC 		& 2:384 \\
					\end{tabular}
					% HDDC = woodford's DC diss; JSR = Marquardt's JS Revelations; TCHC = Vogel HC
				
				\begin{description}
					\item[Online Access] \href{https://www.josephsmithpapers.org/paper-summary/discourse-12-february-1836/1}{Here}
					% \item[Analysis] \hyperref[XX]{Here}
				\end{description}
				
				\begin{description}
					\item[Background]
				\end{description}
				
					% It might also be useful to give a two sentence context of the period: e.g., "This speech was given in Nauvoo to a small congregation one month prior to the destruction of the Nauvoo Expositor. These and other tensions may have been on the prophet's mind when he gave the speech."
				
				\begin{description}
					\item[Original Source]
				\end{description}
				
					% brief description of original source material, with reference to JSP or somewhere else for more info
					% Background material: it might be helpful to have a note that says whether a given document was presented as a speech, was written in a journal, etc.
					% Also a comment here about how reliable the source might be, or what shape it is in, if there are large omissions or parts that are hard to understand, and so on.
					
					% Note the relationship between this entry and the minutes for 12 february; see the journal for that date (J1:185) for the transition.
				
				\begin{description}
					\item[Text]
				\end{description}
				
					I then arose and made some remarks upon the object of our meeting, which were as follows--- first that many are desiring to be ordained to the ministry, who are not called and consequntly the Lord is displeased--- secondly, many already have been ordained who ought [not] to hold official stations in the church because they dishonour themselves and the church and bring persecution swiftly upon us, in consequence of their zeal without k[n]owledge--- I requested the quorum's to take some measures to regulate the same, I proposed some resolutions and remarked to the brethren that the subject was now before them and open for discussion
					
		% -----
		\subsubsection{Letter to Henrietta Raphael Seixas, between 6 and 13 February 1836}
		% -----
			\label{EMS-1836-JS-6-13-FEB-1836}
			% \label{EMS-1830-JS-DAY-3LETTERMONTH-YEAR}
			
			\begin{description}
					\item[Print Sources]
				\end{description}

					\begin{tabular}{p{0.5in}p{1in}p{0.5in}p{1in}}
					JSP	 	& D5:173--178		& JSR 		& --- \\
					DC	 	& ---				& HC 		& --- \\
					HDDC 	& ---				& TCHC 		& --- \\
					\end{tabular}
					% HDDC = woodford's DC diss; JSR = Marquardt's JS Revelations; TCHC = Vogel HC
				
				\begin{description}
					\item[Online Access] \href{https://www.josephsmithpapers.org/paper-summary/letter-to-henrietta-raphael-seixas-between-6-and-13-february-1836/1}{Here}
					% \item[Analysis] \hyperref[XX]{Here}
				\end{description}
				
				\begin{description}
					\item[Background]
				\end{description}
				
					% It might also be useful to give a two sentence context of the period: e.g., "This speech was given in Nauvoo to a small congregation one month prior to the destruction of the Nauvoo Expositor. These and other tensions may have been on the prophet's mind when he gave the speech."
				
				\begin{description}
					\item[Original Source]
				\end{description}
				
					% brief description of original source material, with reference to JSP or somewhere else for more info
					% Background material: it might be helpful to have a note that says whether a given document was presented as a speech, was written in a journal, etc.
					% Also a comment here about how reliable the source might be, or what shape it is in, if there are large omissions or parts that are hard to understand, and so on.
				
				\begin{description}
					\item[Text]
				\end{description}
				
					\begin{tightright}
					Kirtland Ohio Feb. 13\supersu{th}\supers{.} 1836.
					\end{tightright}

						\hypertarget{EMS-1836-JS-6-13-FEB-1836-a}%
					Dear
						\marginnote{a}%
					Madam

					We have the privilege of addressing you a few lines through the kindness of Professor Sexias [Joshua Seixas] who we believe has been sent to this institution through the immediate directions of God to promote the cause of truth and benefit a fallen world. We are in this led to be \sout{believe} thankful to our Redeemer in whose Glorious cause we are engaged as we are anxiously desiring to become acquainted with an individual of virtue \& piety who understood perfectly those languages in which the Scriptures of the Old and New Testaments were originally written as our only object is to do good to lay aside error when we discover it forsake evil and follow righteousness and truly be the better prepared and qualified to render assistance to our fellow men and glorify the name of the Lord: in this our expectations are fully realized and we trust through the goodness of \sout{Good} God to make a proper improvement of the blessing thus given. And we sincerely pray that on the part of your husband our acquaintance may be of that kind that we shall ever have cause to bless and adore God for thus guiding him to this place by his unseen hand.

						\hypertarget{EMS-1836-JS-6-13-FEB-1836-b}%
					We
						\marginnote{b}%
					have seen the possession of Mr. Seixas a very valuable Lexicon which he informs us is your individual property. We have no hesitation in saying that is highly valuable by yourself for the convenience and use of a private family; but we do believe that the wisdom and philanthropy which error inspires the heart of the pure and good will forego for a few months these \uline{special} benefits which may be derived from it for a few for the pleasure of beneftting the many. As the Trustees of this institution we have by the request of others as well as expressing our own individual desires to have it the same taken the liberty to thus tresspass upon your time and patience and request the privilege of purchaseing of you through Professor Seixas this Lexicon. It is unnecessary to wholly unnecessary for us to communicate or attempt to the great worth this Lexicon would be to this Institution in our present and future studies as yourself of the fact and we only say that we hope that God may direct you by his Holy Spirit to do right, and we trust the issue in his hands believing you will call to mind the flood of reproach heaped upon us as a people the few advantages we profess in comparison with other\sout{s} institution in consequence of the same and that all favors thus bestowed will be only appreciated and thankfully remembered.

					\begin{tightcenter}
							\hypertarget{EMS-1836-JS-6-13-FEB-1836-c}%
						We
							\marginnote{c}%
						are most respectfully.

						Your Obt. Servants.

						Joseph Smith Junr.

						Sidney Rigdon

						F[rederick] G. Williams

						Oliver Cowdery
					\end{tightcenter}

					Kirtland Ohio, Feb. 1836

					A true copy from the original.

					\begin{tightright}
					W[arren] F. Cowdery \uline{Rec}.
					\end{tightright}
					
		% -----
		\subsubsection{Record, 17 February 1836}
		% -----
			\label{EMS-1836-JS-17-FEB-1836}
			% \label{EMS-1830-JS-DAY-3LETTERMONTH-YEAR}
			
			\begin{description}
					\item[Print Sources]
				\end{description}

					\begin{tabular}{p{0.5in}p{1in}p{0.5in}p{1in}}
					JSP	 	& ---		& JSR 		& --- \\
					DC	 	& ---		& HC 		& --- \\
					HDDC 	& ---		& TCHC 		& --- \\
					TPJS 	& ---		& 	 		&  \\
					\end{tabular}
					% HDDC = woodford's DC diss; JSR = Marquardt's JS Revelations; TCHC = Vogel HC
				
				\begin{description}
					\item[Online Access] \href{https://www.josephsmithpapers.org/paper-summary/journal-1835-1836/159}{Here}
					% \item[Analysis] \hyperref[XX]{Here}
				\end{description}
				
				\begin{description}
					\item[Background]
				\end{description}
				
					% It might also be useful to give a two sentence context of the period: e.g., "This speech was given in Nauvoo to a small congregation one month prior to the destruction of the Nauvoo Expositor. These and other tensions may have been on the prophet's mind when he gave the speech."
				
				\begin{description}
					\item[Original Source]
				\end{description}
				
					% brief description of original source material, with reference to JSP or somewhere else for more info
					% Background material: it might be helpful to have a note that says whether a given document was presented as a speech, was written in a journal, etc.
					% Also a comment here about how reliable the source might be, or what shape it is in, if there are large omissions or parts that are hard to understand, and so on.
				
				\begin{description}
					\item[Text]
				\end{description}
				
					Wednesday the 17\supers{th} attend[ed] the school and read and translated with my class as usual, and my soul delights in reading the word of the Lord in the original, and I am determined to persue the study of languages untill I shall become master of them, if I am permitted to live long enough, at any rate so long as I do live I am determined to make this my object, and with the blessing of God I shall succe[e]d to my sattisfaction,---
					
		% -----
		\subsubsection{Record, 19 February 1836}
		% -----
			\label{EMS-1836-JS-19-FEB-1836}
			% \label{EMS-1830-JS-DAY-3LETTERMONTH-YEAR}
			
			\begin{description}
					\item[Print Sources]
				\end{description}

					\begin{tabular}{p{0.5in}p{1in}p{0.5in}p{1in}}
					JSP	 	& ---		& JSR 		& --- \\
					DC	 	& ---		& HC 		& --- \\
					HDDC 	& ---		& TCHC 		& --- \\
					TPJS 	& ---		& 	 		&  \\
					\end{tabular}
					% HDDC = woodford's DC diss; JSR = Marquardt's JS Revelations; TCHC = Vogel HC
				
				\begin{description}
					\item[Online Access] \href{https://www.josephsmithpapers.org/paper-summary/journal-1835-1836/159}{Here}
					% \item[Analysis] \hyperref[XX]{Here}
				\end{description}
				
				\begin{description}
					\item[Background]
				\end{description}
				
					% It might also be useful to give a two sentence context of the period: e.g., "This speech was given in Nauvoo to a small congregation one month prior to the destruction of the Nauvoo Expositor. These and other tensions may have been on the prophet's mind when he gave the speech."
				
				\begin{description}
					\item[Original Source]
				\end{description}
				
					% brief description of original source material, with reference to JSP or somewhere else for more info
					% Background material: it might be helpful to have a note that says whether a given document was presented as a speech, was written in a journal, etc.
					% Also a comment here about how reliable the source might be, or what shape it is in, if there are large omissions or parts that are hard to understand, and so on.
				
				\begin{description}
					\item[Text]
				\end{description}
				
					I conversed with Mr. Seixas upon the subject of religion, at my house this afternoon, he listened with attention and appeared interested with my remarks, and I believe the Lord is striving with him, by his holy spirit, and that he will eventually embrace the new and everlasting covenant, for he is a chosen vessel unto the Lord to do his people good,--- but I forbear lest I get to prophesying upon his head--- this evening President Rigdon and myself called at Mr. Seixas lodgings and conversed with him upon the subject of the School, had a pleasant interview
					
		% -----
		\subsubsection{Record, 23 February 1836}
		% -----
			\label{EMS-1836-JS-23-FEB-1836}
			% \label{EMS-1830-JS-DAY-3LETTERMONTH-YEAR}
			
			\begin{description}
					\item[Print Sources]
				\end{description}

					\begin{tabular}{p{0.5in}p{1in}p{0.5in}p{1in}}
					JSP	 	& ---		& JSR 		& --- \\
					DC	 	& ---		& HC 		& --- \\
					HDDC 	& ---		& TCHC 		& --- \\
					TPJS 	& ---		& 	 		&  \\
					\end{tabular}
					% HDDC = woodford's DC diss; JSR = Marquardt's JS Revelations; TCHC = Vogel HC
				
				\begin{description}
					\item[Online Access] \href{https://www.josephsmithpapers.org/paper-summary/journal-1835-1836/162}{Here}
					% \item[Analysis] \hyperref[XX]{Here}
				\end{description}
				
				\begin{description}
					\item[Background]
				\end{description}
				
					% It might also be useful to give a two sentence context of the period: e.g., "This speech was given in Nauvoo to a small congregation one month prior to the destruction of the Nauvoo Expositor. These and other tensions may have been on the prophet's mind when he gave the speech."
				
				\begin{description}
					\item[Original Source]
				\end{description}
				
					% brief description of original source material, with reference to JSP or somewhere else for more info
					% Background material: it might be helpful to have a note that says whether a given document was presented as a speech, was written in a journal, etc.
					% Also a comment here about how reliable the source might be, or what shape it is in, if there are large omissions or parts that are hard to understand, and so on.
				
				\begin{description}
					\item[Text]
				\end{description}
				
					...This afternoon the sisters met again at the chapel to work on the v[e]il toward the close of the day I met with the presidency \& many of the brethren in the house of the Lord--- I made some remarks from the pulpit upon the rise and progress of the church of Christ of Latter day Saints and pronounced a blessing upon the Sisters for the liberality in giving their servises so cheerfully to make the veil for the Lord's house also upon the congregation and dismissed
					
		% -----
		\subsubsection{Record, 3 March 1836}
		% -----
			\label{EMS-1836-JS-3-MAR-1836}
			% \label{EMS-1830-JS-DAY-3LETTERMONTH-YEAR}
			
			\begin{description}
					\item[Print Sources]
				\end{description}

					\begin{tabular}{p{0.5in}p{1in}p{0.5in}p{1in}}
					JSP	 	& ---		& JSR 		& --- \\
					DC	 	& ---		& HC 		& --- \\
					HDDC 	& ---		& TCHC 		& --- \\
					TPJS 	& ---		& 	 		&  \\
					\end{tabular}
					% HDDC = woodford's DC diss; JSR = Marquardt's JS Revelations; TCHC = Vogel HC
				
				\begin{description}
					\item[Online Access] \href{https://www.josephsmithpapers.org/paper-summary/journal-1835-1836/166}{Here}
					% \item[Analysis] \hyperref[XX]{Here}
				\end{description}
				
				\begin{description}
					\item[Background]
				\end{description}
				
					% It might also be useful to give a two sentence context of the period: e.g., "This speech was given in Nauvoo to a small congregation one month prior to the destruction of the Nauvoo Expositor. These and other tensions may have been on the prophet's mind when he gave the speech."
				
				\begin{description}
					\item[Original Source]
				\end{description}
				
					% brief description of original source material, with reference to JSP or somewhere else for more info
					% Background material: it might be helpful to have a note that says whether a given document was presented as a speech, was written in a journal, etc.
					% Also a comment here about how reliable the source might be, or what shape it is in, if there are large omissions or parts that are hard to understand, and so on.
				
				\begin{description}
					\item[Text]
				\end{description}
				
						\hypertarget{EMS-1836-JS-3-MAR-1836-a}%
					...This
						\marginnote{a}%
					evening the several quorums met agreeably to adjourment and were organized according to their official standing in the church.
					
					I then arose and made some remarks on the object of our meeting which are as follows.
					
					1\supers{st} To receive or reject certain resolutions that were drafted by a commitee chosen for that purpose at a preceeding meeting respecting licenses for elders and other official members.
					
					2\supers{nd} To sanction by the united voice of the quorum certain resolutions respecting ordaining members; that had passed through each quorum seperately \sout{for} without any alteration or amendment except<ing> in the quorum of the twelve.
					
						\hypertarget{EMS-1836-JS-3-MAR-1836-b}%
					The
						\marginnote{b}%
					council opened by singing and prayer. President O[liver] Cowdery then arose and read the resolutions respecting licenses three times, the third time he read the \sout{article} <resolutions> he gave time and oppertunity after reading each article for objections to be made if any there were; no objections were made--- I then observed that these resolutions must needs pass through each quorum seperately begining at the presidency and concequently it must first be thrown into the hands of the president of the Deacon<s> \& his council as equal rights \& privileges are my motto, and one man is as good as another, if he behaves as well, and that all men should be esteemed alike, without regard to distinction's of an official nature,---...

					...I then made some remarks on the amendment of the 12. upon the resolutions recorded on pages 155 \& 156...
					
		% -----
		\subsubsection{Record, 12 March 1836}
		% -----
			\label{EMS-1836-JS-12-MAR-1836}
			% \label{EMS-1830-JS-DAY-3LETTERMONTH-YEAR}
			
			\begin{description}
					\item[Print Sources]
				\end{description}

					\begin{tabular}{p{0.5in}p{1in}p{0.5in}p{1in}}
					JSP	 	& ---				& JSR 		& --- \\
					DC	 	& ---				& HC 		& --- \\
					HDDC 	& ---				& TCHC 		& --- \\
					TPJS 	& 108--109			&  &  \\
					\end{tabular}
					% HDDC = woodford's DC diss; JSR = Marquardt's JS Revelations; TCHC = Vogel HC
				
				\begin{description}
					\item[Online Access] \href{https://www.josephsmithpapers.org/paper-summary/journal-1835-1836/170}{Here}
					% \item[Analysis] \hyperref[XX]{Here}
				\end{description}
				
				\begin{description}
					\item[Background]
				\end{description}
				
					% It might also be useful to give a two sentence context of the period: e.g., "This speech was given in Nauvoo to a small congregation one month prior to the destruction of the Nauvoo Expositor. These and other tensions may have been on the prophet's mind when he gave the speech."
				
				\begin{description}
					\item[Original Source]
				\end{description}
				
					% brief description of original source material, with reference to JSP or somewhere else for more info
					% Background material: it might be helpful to have a note that says whether a given document was presented as a speech, was written in a journal, etc.
					% Also a comment here about how reliable the source might be, or what shape it is in, if there are large omissions or parts that are hard to understand, and so on.
				
				\begin{description}
					\item[Text]
				\end{description}
				
					I was informed to day that a man by the name of Clark froze to death, last night near this place, who was under the influenc of ardent spirits; O my God how long will this monster intemperance find it's victims on the earth, me thinks until the earth is swept with the wrath \sout{of} and indignation of God, and christ's kingdom becomes universal. O come Lord Jesus and cut short thy work in rightieousness.
					
		% -----
		\subsubsection{Record, 19 March 1836}
		% -----
			\label{EMS-1836-JS-19-MAR-1836}
			% \label{EMS-1830-JS-DAY-3LETTERMONTH-YEAR}
			
			\begin{description}
					\item[Print Sources]
				\end{description}

					\begin{tabular}{p{0.5in}p{1in}p{0.5in}p{1in}}
					JSP	 	& ---				& JSR 		& --- \\
					DC	 	& ---				& HC 		& --- \\
					HDDC 	& ---				& TCHC 		& --- \\
					\end{tabular}
					% HDDC = woodford's DC diss; JSR = Marquardt's JS Revelations; TCHC = Vogel HC
				
				\begin{description}
					\item[Online Access] None known %\href{https://www.josephsmithpapers.org/paper-summary/journal-1835-1836/170}{Here}
					% \item[Analysis] \hyperref[XX]{Here}
				\end{description}
				
				\begin{description}
					\item[Background]
				\end{description}
				
					% It might also be useful to give a two sentence context of the period: e.g., "This speech was given in Nauvoo to a small congregation one month prior to the destruction of the Nauvoo Expositor. These and other tensions may have been on the prophet's mind when he gave the speech."
				
				\begin{description}
					\item[Original Source]
				\end{description}
				
					% brief description of original source material, with reference to JSP or somewhere else for more info
					% Background material: it might be helpful to have a note that says whether a given document was presented as a speech, was written in a journal, etc.
					% Also a comment here about how reliable the source might be, or what shape it is in, if there are large omissions or parts that are hard to understand, and so on.
					
					This is in the Kirtland EQ Record, edited by Cook and Backman, page 12.
				
				\begin{description}
					\item[Text]
				\end{description}
				
					This evening the Elder[s] met at the Lord's house, with other quorums of the church. The meeting was opened by singing, and prayer. President Joseph Smith J[u]n. gave some instructions relative to the sealing power, in a short but powerful address.
					
					President S. Rigdon sealed the blessings of the Lord on those who had been anointed, by prayer and with a shout of Hosanna to God and, the Lamb.
					
					Some of the brethren spake of the goodness of the Lord, and President J. Smith jun made some remarks respecting the coming meeting, and gave out some appointments...
					
		% -----
		\subsubsection{Record, 27 March 1836}
		% -----
			\label{EMS-1836-JS-27-MAR-1836-REC}
			% \label{EMS-1830-JS-DAY-3LETTERMONTH-YEAR}
			
			\begin{description}
					\item[Print Sources]
				\end{description}

					\begin{tabular}{p{0.5in}p{1in}p{0.5in}p{1in}}
					JSP	 	& ---				& JSR 		& --- \\
					DC	 	& ---				& HC 		& --- \\
					HDDC 	& ---				& TCHC 		& --- \\
					\end{tabular}
					% HDDC = woodford's DC diss; JSR = Marquardt's JS Revelations; TCHC = Vogel HC
				
				\begin{description}
					\item[Online Access] \href{https://www.josephsmithpapers.org/paper-summary/journal-1835-1836/174}{Here}
					% \item[Analysis] \hyperref[XX]{Here}
				\end{description}
				
				\begin{description}
					\item[Background]
				\end{description}
				
					% It might also be useful to give a two sentence context of the period: e.g., "This speech was given in Nauvoo to a small congregation one month prior to the destruction of the Nauvoo Expositor. These and other tensions may have been on the prophet's mind when he gave the speech."
				
				\begin{description}
					\item[Original Source]
				\end{description}
				
					% brief description of original source material, with reference to JSP or somewhere else for more info
					% Background material: it might be helpful to have a note that says whether a given document was presented as a speech, was written in a journal, etc.
					% Also a comment here about how reliable the source might be, or what shape it is in, if there are large omissions or parts that are hard to understand, and so on.
				
				\begin{description}
					\item[Text]
				\end{description}
				
						\hypertarget{EMS-1836-JS-27-MAR-1836-REC-a}%
					Sunday
						\marginnote{a}%
					morning the 27\supers{th} The congregation began to \sout{assembly} <assimble> <at the chapel> at about 7 oclock one hour \sout{early} <earlier> than the doors were to be opened many brethren had come in from the region's round about to witness the dedication of the Lord's House and share in his blessings and such was the anxiety on this occasion that some hundreds, (probably five or six,) assembled \sout{collected} before the doors were opened---

						\hypertarget{EMS-1836-JS-27-MAR-1836-REC-b}%
					The
						\marginnote{b}%
					presidency entered with the door kepers and aranged them at the inner and outer doors also placed our stewards to receiv donations from those who should feel disposed to contribute something to defray the expenses of building the House of the Lord--- <we also dedicated the pulpits \& consecrated them to the Lord> The doors were then opened President [Sidney] Rigdon President [Oliver] Cowdery and myself seated the congregation as they came in, \sout{we} \sout{received} \sout{about} and according to the best calculation we could make we received between 9 and 10,00 <hundred> which \sout{was} is as many as can be comfortably situated we then informed the door keepers that we could receive no more, and a multitude were deprived of the benefits of the meeting on account of the house not being sufficiently capacious to receive them, and I felt to regret that any of my brethren and sisters should be deprived of the Meeting, and I recommended them to repair to the School-house and hold a meeting which they did and filled that house <also> and yet many were left out---

						\hypertarget{EMS-1836-JS-27-MAR-1836-REC-c}%
					The
						\marginnote{c}%
					assembly were then organized in the following manner,--- viz.

					\sout{President F G. Williams Presdt. Joseph Smith<, Sen.> and Presdt W W. Phelps occupied the first pulpit <in the west end of the house> for the Melchisedec priesthood, Presdt. S. Rigdon, Myself, and Prest. Hyrum Smith the 2\supers{ond}--- Presdt D. Whitmer Presdt O. Cowdery and Presdt John. Whitmer the 3\supers{d}--- The 4\supers{th} was occupied by the president of the high priests and his counsellors, and 2 choirister's--- The 12. apostles on the right, the high council of Kirtland on the left, The pulpits <in the east end of the house> for the Aaronic priesthood were occupied in the following manner--- The Bishop of Kirtland and his counsellors in the first pulpit--- <the> Bishop of Zion and his counsellors in the 2\supers{ond}--- the presdt. of the priest and his counsellors in the 3\supers{d}--- the presdt. of the Teachers and his counsellors in the 4\supers{th}--- the high council of Zion on the right, the 7. presdt of the Seventies on the left}

					\begin{tightcenter}
							\hypertarget{EMS-1836-JS-27-MAR-1836-REC-d}%
						West
							\marginnote{d}%
						end of the. house
					\end{tightcenter}

					Presdt. F[rederick] G. Williams Presdt. Joseph Smith, Sen and Presdt. W[illiam] W. Phelps occupied the 1\supers{st} pulpit for the Melchisedic priesthood--- Presdt. S. Rigdon myself and Presdt Hyrum Smith in the 2\supers{ond}---

					Presdt. D[avid] Whitmer Presdt. O. Cowdery and Presdt. J[ohn] Whitmer in the 3\supers{d.}--- The 4th was occupied by the president of the high-preists and his counsellors, and 2 choiristers--- The 12. Apostles on the right in the 3. highest seats---

					The presdt of the Eldrs his clerk \& counsellors in the seat immediatly below the 12--- The high council of Kirtland consisting of 12, on the left in the 3, first seats--- the 4\supers{th} seat below them was occupied by \sout{the} \sout{presidency's} Eldr's W[arren] A. Cowdery \& W[arren] Parrish who served as scribes.---

						\hypertarget{EMS-1836-JS-27-MAR-1836-REC-e}%
					The
						\marginnote{e}%
					pulpits in the east end of the house for the Aaronic priesthood were occupied as follows.--- The Bishop of Kirtland and his counsellors in the 1\supers{st} pulpit,--- The Bishop of Zion and his counsellors in the 2\supers{ond}--- The presdt. of the priests and his counsellors in the 3\supers{d.}--- The presdt. of the Teachers \sout{in} and his counsellors <\& one choirister> in the 4\supers{th}--- The high council of Zion consisting of 12. counsellors on the right--- The presdt of the Deacons and his counsillors in the seat below them.--- The 7. presdts of the Seventies on the left--- The choir of singers were seated in the 4 corners of the room in seats prepared for that purpose--- <rec\supers{d} by contribution \$960.00>

					9 oclock A. M the servises of the day were opened by Presdt S. Rigdon by reading 1\supers{st} the 96 Psalm secondly the 24\supers{th} Psalm--- the choir then sung hymn on the 29\supers{th} page of Latter day Saints collection of hymn's--- prayer by Presdt Rigdon choir then sung hymn on 14\supers{th} page Presdt Rigdon then <read> the 18, 19, \& 20, verses of the 8\supers{th} chapter of Mathew and preached more particularly from the 20\supers{th} verse.--- his prayer and address were very \sout{forcibly} <forcible> and sublime, and well adapted to the occasion.--- after he closed his \sout{services} <sermon>, he called upon the several quorums commenceing with the presidency, to manifest by rising up, their willingness to acknowledge me as a prophet and seer and uphold me as such by their p[r]ayers of faith, all the quorums in their turn, cheerfully complyed with this request he then called upon all the congregation of Saints, also to give their assent by rising on their feet which they did unanimously

						\hypertarget{EMS-1836-JS-27-MAR-1836-REC-f}%
					After
						\marginnote{f}%
					an intermission of 20, minutes the servises of the day were resumed, by singing Adam ondi ahman. I then made a short address and called upon the several quorums, and all the congregation of saints to acknowledge the Presidency as Prophets and Seers, and uphold them by their prayers, they all covenanted to do so by rising; I then called upon the quorums and congregation of saints to acknowledge the 12 Apostles who were present as Prophets and Seers and special witnesses to all the nations of the earth, holding the keys of the kingdom, to unlock it or cause it to be done among \sout{all} \sout{nations} them; and uphold them by their prayers, which they assented to by rising,

						\hypertarget{EMS-1836-JS-27-MAR-1836-REC-g}%
					I
						\marginnote{g}%
					then called upon the quorums and congregation of saints to acknowledge the high council of Kirtland in all the authorities [authority] of the Melchisedec priesthood and uphold them by their prayers which they assented to by rising. I then called upon the quorums and congregation of saints to acknowledge and uphold by their prayer's the Bishops of Kirtland and Zion and their counsellors, \sout{the} \sout{Presidents} \sout{of} \sout{the} \sout{Priests} in all the authority of the Aaronic priesthood, which they did by rising. I then called upon the quorums and congregation of saints to acknowledge the high-council of Zion, and uphold them by their prayers in all the authority of the high priesthood which they did by rising. I next called upon the quorums and congregation of saints to acknowledge the Presidents of the seventy's who act as their representives as <Apostles and> special witnesses to the nations to assist the 12 in opening the gospel kingdom, among all people and to uphold them by their prayer's which they did by rising--- I then called upon the quorums and all the saints to acknowledge president of the Elders and his counsellors and uphold them by their prayers which they did by rising---. The quorums and congregation of saints were then called upon to acknowledge and uphold by their prayer's the Presidents of the Priests, Teachers, and Deacons and their counsellors, which they did by rising.

					N. B. The Presidents \sout{were} of the seventy's were acknowledged first after the 12 Apostles

					The hymn on the hundred and 14 page was then sung, after which I offered to God the following dedication prayer.

					\begin{tightcenter}
							\hypertarget{EMS-1836-JS-27-MAR-1836-REC-h}%
						Prayer,
							\marginnote{h}
					\end{tightcenter}

					At the dedication of the Lord's House in Kirtland Ohio March 27, 1836.--- by Joseph Smith, jr. President of the Church of the Latter Day Saints.

					Thanks be to thy name, O Lord God of Israel, who keepest covenant and shewest mercy unto thy servants, who walk uprightly before thee with all their hearts; thou who hast commanded thy servants to build an house to thy name in this plase. (Kirtland.) And now thou beholdest, O Lord, that so thy servants have done, according to thy commandment. And now we ask the[e], holy Father, in the name of Jesus Christ, the Son of thy bosom, in in whose name alone salvation can be administered to the children of men: we ask the[e], O Lord, to accept of this house, the workmanship of the hands of us, thy servants, which thou didst command us to build; for thou knowest that we have done this work through great tribulation: and out of our poverty we have given of our substance to build a house to thy name, that the Son of Man might have a place to manifest himself to his people.

						\hypertarget{EMS-1836-JS-27-MAR-1836-REC-i}%
					And
						\marginnote{i}%
					as thou hast said, in a revelation given unto us, calling us thy friends, saying--- ``Call your solemn assembly, as I have commanded you; and as all have not faith, seek ye diligently and teach one another words of wisdom; yea, seek ye out of the best books words of wisdom: Seek learning, even by study, and also by faith.

					``Organize yourselves; prepare every \sout{thing} needful thing, and establish a house, even a house of prayer, a house [of] fasting, a house of faith, a house of learning

					a house of glory, a house of order, a house of God: that your incomings may be in the name of the Lord, that your out goings may be in the name of the Lord: that all your salutations may be in the name of the Lord, with uplifted hands to the Most High.''

						\hypertarget{EMS-1836-JS-27-MAR-1836-REC-j}%
					And
						\marginnote{j}%
					now, Holy Father, we ask thee to assist us, thy people with thy grace in calling our solemn assembly, that it may be done to thy honor, and to thy divine acceptance, and in a manner that we may be found worthy in thy sight, to secure a fulfilment of the promises which thou hast made unto us thy people, in the revelatio[n]s given unto us: that thy glory may rest down upon thy people, and upon this thy house, which we now dedicate to thee, that it may be sanctified and consecrated to be holy, and that thy holy presence may be continually in this house; and that all people who shall enter upon the threshold of the Lord's house may feel thy power and be constrained to acknowledge that thou hast sanctified it, and that it is thy house, a place of thy holiness.

						\hypertarget{EMS-1836-JS-27-MAR-1836-REC-k}%
					And
						\marginnote{k}%
					do thou grant, holy Father, that all those who shall worship in this house, may be taught words of wisdom out of the best books, and that they may seek learning, even by study, and also by faith; as thou hast said; and that they may grow up in thee and receive a fulness of the Holy Ghost, and be organized according to thy laws, and be prepared to obtain every needful thing, and that this house may be a house of prayer, a house of fasting, a house of faith, a house of glory, and of God, even thy house: that all the incomings of thy people, into this house, may be in the name of the Lord; that all their outgoings, from this house, may be in the name of the Lord; that all their salutations may be in the name of

					Lord, with holy hands uplifted to the Most High; and that no unclean thing shall be permitted to come into thy house to pollute it.

					And when thy people transgress, any of them, they may speedily repent and return unto thee, and find favour in thy sight, and be restored to the blessings which thou hast ordained, to be poured out upon those who shall reverence thee in this thy house.

						\hypertarget{EMS-1836-JS-27-MAR-1836-REC-l}%
					And
						\marginnote{l}%
					we ask, holy Father, that thy servants may go forth from this house, armed with thy power, and that thy name may be upon them and thy glory be round about them, and thine angels have charge over them; and from this place they may bear exceeding great and glorious tidings, in truth, unto the ends of the earth, that they may know that this is thy work, and that thou hast put forth thy hand, to fulfil that which thou hast spoken by the mouths of thy prophets concerning the last days.

					We ask the[e]; holy Father; to establish the people that shall worship and honorably hold a name and standing in this thy house, to all generations, and for eternity that no weapon formed against them shall prosper; that he who diggeth a pit for them shall fall into the same himself; that no combination of wickedness shall have power to rise up and prevail over thy people, upon whom thy name shall be put in this house: and if any people shall rise against this people, that thine anger be kindled against them: and if they shall smite this people, thou wilt smite them--- thou wilt fight for thy people as thou didst in the day of battle, that they may be delivered from the hands of all their enimies.

						\hypertarget{EMS-1836-JS-27-MAR-1836-REC-m}%
					We
						\marginnote{m}%
					ask thee, holy Father, to confound, and astonish, and bring to shame, and confusion, all those who have

					spread lying reports abroad over the world against thy servant or servants, if they will not repent when the everlasting gospel shall be proclaimed in their ears, and that all their works may be brought to nought, and be swept away by the hail, and by the judgments, which thou wilt send upon them in thine anger, that their may be an end to lyings and slanders against thy people: for thou knowest, O Lord, that thy servants have been innocent before thee in bearing record of thy name for which they have suffered these things; therefore we plead before thee for a full and complete deliverence from under this yoke. Break it off O Lord: break it off from the necks of thy servants, by thy power, that we may rise up in the midst of this generation and do thy work!

						\hypertarget{EMS-1836-JS-27-MAR-1836-REC-n}%
					O
						\marginnote{n}%
					Jehovah, have mercy upon this people, and as all men sin, forgive the transgressions of thy people, and let them be blotted out forever. Let the anointing of thy ministers be sealed upon them with power from on high: let it be fulfilled upon them as upon those on the day of Pentacost: let the gift of tongues be poured out upon thy people, even cloven tongues as of fire, and the interpretation thereof. And let thy house be filled, as with a rushing mighty wind, with thy glory.

					Put upon thy servants the testimony of the covenant that where they go out and proclaim thy word, they may seal up the law, and prepare the hearts of thy saints for all those judgements thou art about to send, in thy wrath, upon the inhabitants of the earth because of their transgressions, that thy people may not faint in the day of trouble.

						\hypertarget{EMS-1836-JS-27-MAR-1836-REC-o}%
					And
						\marginnote{o}%
					whatever city thy servants shall enter, and the people of that city receive their testimony, let thy peace and thy salvation be upon that city, that they may gather out from that city the righteous, that they may come forth to Zion, or to her stakes, the places of thine appointment, with songs of everlasting joy,--- and until this be accomplished let not thy judgements fall upon that city.

					And whatever city thy servants shall enter, and the people of that city receive not the\sout{ir} testimony of thy servants, and thy servants warn them to save themselves from this untoward generation let it be upon that city according to that which thou hast spoken, by the mouths of thy prophets; but deliver thou, O Jehovah, we beseech thee, thy servants from their hands, and cleanse them from their blood. O Lord, we delight not in the destruction of our fellow men: their souls are precious before thee; but thy word must be fulfilled:--- help thy servants to say, with thy grace assisting them, thy will be done, O Lord, and not our<s>.

						\hypertarget{EMS-1836-JS-27-MAR-1836-REC-p}%
					We
						\marginnote{p}%
					know that thou hast spoken by the mouth of thy prophets, terrible things concerning the wicked in the last days, that thou wilt pour out thy judgements; without measure; therefore, O Lord, deliver thy people from the calamity of the wicked; enable thy servants to seal up the law and bind up the testimony, that they may be prepared against the day of burning.

					We ask thee, holy Father, to remember those who have been driven by the inhabitants of Jackson county Missouri, from the lands of their inheritance, and break off, O Lord, this yoke of affliction

					that has been put upon them. Thou knowest, O Lord, that they have been greatly oppressed and afflicted, by wicked men, and our hearts flow out in sorrow because of their grevious burdens. O Lord, how long wilt thou suffer this people to bear this affliction, and the cries of the innocent ones to ascend up in thine ears, and their blood to come up in testimony before thee and not make a display of thy power in their behalf?

						\hypertarget{EMS-1836-JS-27-MAR-1836-REC-q}%
					Have
						\marginnote{q}%
					mercy, O Lord, upon that wicked mob, who have driven thy people, that they may cease to spoil, that they may repent of their sins, if repentance is to be found; but if they will not, make bare thine arm, O Lord, and redeem that which thou didst appoint a Zion unto thy people.

					And if it cannot be otherwise, that the cause of thy people may not fail before thee, may thine anger be kindled and thine indignation fall upon them that they may be wasted away, both root and branch from under heaven; but inasmuch as they will repent, thou art gracious and merciful, and will turn away thy wrath, when thou lookest upon the face of thine anointed.

					Have mercy, O Lord, upon all the nations of the earth: have mercy upon the rulers of our land may those principles which were so honorably and nobly defended: viz, the constitution of our land, by our fathers, be established forever.

						\hypertarget{EMS-1836-JS-27-MAR-1836-REC-r}%
					Remember
						\marginnote{r}%
					the kings, the princes, the nobles, and the great ones of the earth, and all people; and the churches: all the poor, the needy and the afflicted ones of the earth, that their hearts may be softened when thy servants shall go out from thy house, O

					Jehovah, to bear testimony of thy name, <that> their prejudices may give way before the truth, and thy people may obtain favour in the sight of all, that all the ends of the earth may know that we thy servants have heard thy voice, and that thou hast sent us, that from among all these thy servants, the sons of Jacob, may gather out the righteous to build a holy city to thy name, as thou hast commanded them.

					We ask thee to appoint unto Zion other stakes besides this one, which thou hast appointed, that the gathering of thy people may roll on in great power and majesty, that thy work may be cut short in righteousness.

						\hypertarget{EMS-1836-JS-27-MAR-1836-REC-s}%
					Now
						\marginnote{s}%
					these words, O Lord, we have spoken before thee, concerning the revelations and commandments which thou hast given unto us, who are i[de]ntified with the Gentiles;--- But thou knowest that we have a great love for the children of Jacob who have been scattered upon the mountains; for a long time in a cloudy and dark day.

					We therefore ask thee to have mercy upon the children of Jacob, that Jerusalem, from this hour, may begin to be redeemed; and the yoke of bondage may begin to be broken off from the house of David, and the children of Judah may begin to return to the lands which thou didst give to Abraham, their father, and cause that the remnants of Jacob, who have been cursed and smitten, because of their transgression, to be converted from their wild and savage condition, to the fulness of the everlasting gospel, that they may lay down their weapons of bloodshed and cease their rebellions. And may
					
					all the scattered remnants of Israel, who have been driven to the ends of the earth, come to a knowledge of the truth, believe in the Messiah, and be redeemed from oppression, and rejoice before thee.

						\hypertarget{EMS-1836-JS-27-MAR-1836-REC-t}%
					O
						\marginnote{t}%
					Lord, remember thy servant Joseph Smith jr. and all his afflictions and persecutions, how he has covenanted with Jehovah and vowed to thee O mighty God of Jacob, and the commandments which thou hast given unto him, and that he hath sincerely strove to do thy will.--- Have mercy, O Lord, upon his wife and children, that they may be exalted in thy presence, and preserved by thy fostering hand,--- Have mercy upon all their immediate connexions, that their prejudices may be broken up, and swept away as with a flood, that they may be converted and redeemed with Israel and know that thou art God.

					Remember, O Lord, the presidents, even all the presidents of thy church, that thy right hand may exalt them with all their families, and their immediate connexions, that their names may be perpetuated and had in everlasting remembrance from generation to generation.

					Remember all thy church, O Lord, with all their families, and all their immediate connexions, with all their sick and afflicted ones, with all the poor and meek of the earth; that the kingdom which thou hast set up without hands, may become a great mountain and fill the whole earth, that thy church may come forth out of the wilderness of darkness, and shine forth fair as the moon, clear as the sun, and terrible as an army with banners, and be adorned as a bride for that day when

						\hypertarget{EMS-1836-JS-27-MAR-1836-REC-u}%
					thou
						\marginnote{u}%
					shalt unveil the heavens, and cause the mountains to flow down at thy presence, and the valley's to be exalted, the rough places made smooth; that thy glory may fill the earth.

					That when the trump shall sound for the dead, we shall be caught up in the cloud to meet thee, that we may ever be with the Lord, that our garments may be pure, that we may be clothed upon with robes of righteousness, with palms in our hands, and crowns of glory upon our heads, and reap eternal joy for all our sufferings.

					O Lord, God Almighty, hear us in these our petitions, and answer us from heaven, thy holy habitation, where thou sittest enthroned, with glory, honour, power majesty, might, dominion, truth, justice judgement, mercy and an infinity of fulness, from everlasting to everlasting.

					O hear, O hear, O hear us, O Lord, and answer these petitions, and accept the dedication of this house, unto thee, the work of our hands, which we have built unto thy name; and also this church to put upon it thy name. And help us by the power of thy spirit, that we may mingle our voices with those bright shining seraphs, around thy throne with acclamations of praise, singing hosanna to God and the Lamb: and let these thine anointed ones be clothed with salvation, and thy saints shout aloud for joy.

					Amen and Amen.

						\hypertarget{EMS-1836-JS-27-MAR-1836-REC-v}%
					Sung
						\marginnote{v}%
					Hosanah to God and the Lamb after which the Lords supper was administered

					I then bore testimony of the administering of angels.--- Presdt Williams also arose and testified that while Presdt Rigdon was making his first prayer an angel entered the window and <took his> seated \sout{himself} between father Smith, and himself, and remained their during his prayer Presdt David Whitmer also saw angels in the house

					We then sealed the proceedings of the day by \sout{a} shouting hosanah to God and the Lamb 3 times sealing it each time with Amen, Amen, and Amen  and after requesting all the official members to meet again in the evening we retired---

					met in the evening and instructed the quorums respecting the ordinance of washing of feet which we were to attend to on wednesday following
					
		% -----
		\subsubsection{Prayer of Dedication, 27 March 1836 [D\&C 109]}
		% -----
			\label{EMS-1836-JS-27-MAR-1836}
			% \label{EMS-1830-JS-DAY-3LETTERMONTH-YEAR}
			
			\begin{description}
					\item[Print Sources]
				\end{description}

					\begin{tabular}{p{0.5in}p{1in}p{0.5in}p{1in}}
					JSP	 	& D5:188--209		& JSR 		& --- \\
					DC	 	& Section 109		& HC 		& 2:410--428 \\
					HDDC 	& 1440--1457		& TCHC 		& 2:399--414 \\
					\end{tabular}
					% HDDC = woodford's DC diss; JSR = Marquardt's JS Revelations; TCHC = Vogel HC
				
				\begin{description}
					\item[Online Access] \href{https://www.josephsmithpapers.org/paper-summary/prayer-of-dedication-27-march-1836-dc-109/1}{Here}
					% \item[Analysis] \hyperref[DC109]{Here}
				\end{description}
				
				\begin{description}
					\item[Background]
				\end{description}
				
					% It might also be useful to give a two sentence context of the period: e.g., "This speech was given in Nauvoo to a small congregation one month prior to the destruction of the Nauvoo Expositor. These and other tensions may have been on the prophet's mind when he gave the speech."
				
				\begin{description}
					\item[Original Source]
				\end{description}
				
					% brief description of original source material, with reference to JSP or somewhere else for more info
					% Background material: it might be helpful to have a note that says whether a given document was presented as a speech, was written in a journal, etc.
					% Also a comment here about how reliable the source might be, or what shape it is in, if there are large omissions or parts that are hard to understand, and so on.
				
				\begin{description}
					\item[Text]
				\end{description}
				
					\begin{tightcenter}
					PRAYER,

					\textit{At the Dedication of the Lord's House in Kirtland, Ohio, March} 27, 1836\textit{,---By JOSEPH SMITH, jr.}

					President of the Church of the Latter Day Saints.
					\end{tightcenter}

					\textsc{Thanks} be to thy name, O Lord God of Israel, who keepest covenant and shewest mercy unto thy servants, who walk uprightly before thee with all their hearts: thou who hast commanded thy servants to build an house to thy name in this place. (Kirtland.) And now thou beholdest, O Lord, that so thy servants have done, according to thy commandment. And now we ask thee, holy Father, in the name of Jesus Christ, the Son of thy bosom, in whose name alone salvation can be administered to the children of men: we ask thee, O Lord, to accept of this house, the workmanship of the hands of us, thy servants, which thou didst command us to build; for thou knowest that we have done this work through great tribulation: and out of our poverty we have given of our substance to build a house to thy name, that the Son of Man might have a place to manifest himself to his people.

					And as thou hast said, in a revelation given unto us, calling us thy friends, saying---``Call your solemn assembly, as I have commanded you; and as all have not faith, seek ye diligently and teach one another words of wisdom; yea, seek ye out of the best books words of wisdom: Seek learning, even by study, and also by faith.

					``Organize yourselves; prepare every needful thing, and establish a house, even a house of prayer, a house of fasting, a house of faith, a house of learning, a house of glory, a house of order, a house of God: that your incomings may be in the name of the Lord; that your out goings may be in the name of the Lord: that all your salutations may be in the name of the Lord, with uplifted hands to the Most High.''

					And now, holy Father, we ask thee to assist us, thy people with thy grace in calling our solemn assembly, that it may be done to thy honor, and to thy divine acceptance, and in a manner that we may be found worthy, in thy sight, to secure a fulfilment of the promises which thou hast made unto us thy people, in the revelations given unto us: that thy glory may rest down upon thy people, and upon this thy house, which we now dedicate to thee, that it may be sanctified and consecrated to be holy, and that thy holy presence may be continually in this house; and that all people who shall enter upon the threshhold of the Lord's house may feel thy power and be constrained to acknowledge that thou hast sanctified it, and that it is thy house, a place of thy holiness.

					And do thou grant, holy Father, that all those who shall worship in this house, may be taught words of wisdom out of the best books, and that they may seek learning, even by study, and also by faith; as thou hast said; and that they may grow up in thee and receive a fulness of the Holy Ghost, and be organized according to thy laws, and be prepared to obtain every needful thing: and that this house may be a house of prayer, a house of fasting, a house of faith, a house of glory, and of God, even thy house: that all the incomings of thy people, into this house, may be in the name of the Lord; that all their outgoings, from this house, may be in the name of the Lord; that all their salutations may be in the name of the Lord, with holy hands, uplifted to the Most High; and that no unclean thing shall be permitted to come into thy house to pollute it.

					And when thy people transgress, any of them, they may speedily repent and return unto thee, and find favor in thy sight, and be restored to the blessings which thou hast ordained, to be poured out upon those who shall reverence thee in this thy house.

					And we ask thee, holy Father, that thy servants may go forth from this house, armed with thy power, and that thy name may be upon them and thy glory be round about them, and thine angels have charge over them; and from this place they may bear exceeding great and glorious tidings, in truth, unto the ends of the earth, that they may know that this is thy work, and that thou hast put forth thy hand, to fulfil that which thou has spoken by the mouths of thy prophets concerning the last days.

					We ask thee, holy Father, to establish the people that shall worship and honorably hold a name and standing in this thy house, to all generations, and for eternity, that no weapon formed against them shall prosper; that he who diggeth a pit for them shall fall into the same himself; that no combination of wickedness shall have power to rise up and prevail over thy people, upon whom thy name shall be put in this house: and if any people shall rise against this people, that thine anger be kindled against them: and if they shall smite this people, thou wilt smite them---thou wilt fight for thy people as thou didst in the day of battle, that they may be delivered from the hands of all their enemies.

					We ask thee, holy Father, to confound, and astonish, and bring to shame, and confusion, all those who have spread lying reports abroad over the world against thy servant, or servants, if they will not repent when the everlasting gospel shall be proclaimed in their ears, and that all their works may be brought to nought, and be swept away by the hail, and by the judgments, which thou wilt send upon them in thine anger, that there may be an end to lyings and slanders against thy people: for thou knowest, O Lord, that thy servants have been innocent before thee in bearing record of thy name for which they have suffered these things; therefore we plead before thee for a full and complete deliverance from under this yoke. Break it off O Lord: break it off from the necks of thy servants, by thy power, that we may rise up in the midst of this generation and do thy work!

					O Jehovah, have mercy upon this people, and as all men sin, forgive the transgressions of thy people, and let them be blotted out forever. Let the annointing of thy ministers be sealed upon them with power from on high: let it be fulfilled upon them as upon those on the day of Pentacost: let the gift of tongues be poured out upon thy people, even cloven tongues as of fire, and the interpretation thereof. And let thy house be filled, as with a rushing mighty wind, with thy glory.

					Put upon thy servants the testimony of the covenant, that when they go out and proclaim thy word, they may seal up the law, and prepare the hearts of thy saints for all those judgements thou art about to send, in thy wrath, upon the inhabitants of the earth, because of their transgressions, that thy people may not faint in the day of trouble.

					And whatever city thy servants shall enter, and the people of that city receive their testimony, let thy peace and thy salvation be upon that city, that they may gather out of that city the righteous, that they may come forth to Zion, or to her stakes, the places of thine appointment, with songs of everlasting joy,---and until this be accomplished let not thy judgments fall upon that city.

					And whatever city thy servants shall enter, and the people of that city receive not the testimony of thy servants, and thy servants warn them to save themselves from this untoward generation, let it be upon that city according to that which thou hast spoken, by the mouths of thy prophets; but deliver thou, O Jehovah, we beseech thee, thy servants from their hands, and cleanse them from their blood. O Lord, we delight not in the destruction of our fellow men: their souls are precious before thee; but thy word must be fulfilled:---help thy servants to say, with thy grace assisting them, thy will be done, O Lord, and not ours.

					We know that thou hast spoken by the mouth of thy prophets, terrible things concerning the wicked, in the last days, that thou wilt pour out thy judgments, without measure: therefore, O Lord, deliver thy people from the calamity of the wicked; enable thy servants to seal up the law and bind up the testimony, that they may be prepared against the day of burning.

					We ask thee, holy Father, to remember those who have been driven by the inhabitants of Jackson county, Missouri, from the lands of their inheritance, and break off, O Lord, this yoke of affliction, that has been put upon them. Thou knowest, O Lord, that they have been greatly oppressed, and afflicted, by wicked men, and our hearts flow out in sorrow because of their grievous burdens. O Lord, how long wilt thou suffer this people to bear this affliction, and the cries of their innocent ones to ascend up in thine ears, and their blood to come up in testimony before thee, and not make a display of thy power in their behalf?

					Have mercy, O Lord, upon that wicked mob, who have driven thy people, that they may cease to spoil, that they may repent of their sins, if repentance is to be found; but if they will not, make bear thine arm O Lord, and redeem that which thou didst appoint a Zion unto thy people!

					And if it can not be otherwise, that the cause of thy people may not fail before thee, may thine anger be kindled and thine indignation fall upon them, that they may be wasted away, both root and branch from under heaven; but in as much as they will repent, thou art gracious and merciful, and will turn away thy wrath, when thou lookest upon the face of thine annointed.

					Have mercy, O Lord, upon all the nations of the earth: have mercy upon the rulers of our land: may those principles which were so honorably and nobly defended: viz, the constitution of our land, by our fathers, be established forever. Remember the kings, the princes, the nobles, and the great ones of the earth, and all people; and the churches: all the poor, the needy and the afflicted ones of the earth, that their hearts may be softened when thy servants shall go out from thy house, O Jehovah, to bear testimony of thy name, that their prejudices may give way before the truth, and thy people may obtain favor in the sight of all, that all the ends of the earth may know that we thy servants have heard thy voice, and that thou hast sent us, that from among all these thy servants, the sons of Jacob, may gather out the righteous to build a holy city to thy name, as thou hast commanded them.

					We ask thee to appoint unto Zion other stakes besides this one, which thou hast appointed, that the gathering of thy people may roll on in great power and majesty, that thy work may be cut short in righteousness.

					Now these words, O Lord, we have spoken before thee, concerning the revelations and commandments which thou hast given unto us, who are identified with the Gentiles;---But thou knowest that we have a great love for the children of Jacob who have been scattered upon the mountains; for a long time in a cloudy and dark day.

					We therefore ask thee to have mercy upon the children of Jacob, that Jerusalem, from this hour, may begin to be redeemed; and the yoke of bondage may begin to be broken off from the house of David, and the children of Judah may begin to return to the lands which thou didst give to Abraham, their father, and cause that the remnants of Jacob, who have been cursed and smitten, because of their transgression, to be converted from their wild and savage condition, to the fulness of the everlasting gospel, that they may lay down their weapons of bloodshed and cease their rebellions. And may all the scattered remnants of Israel, who have been driven to the ends of the earth, come to a knowledge of the truth, believe in the Messiah, and be redeemed from oppression, and rejoice before thee.

					O Lord, remember thy servant Joseph Smith, jr. and all his afflictions and persecutions, how he has covenanted with Jehovah and vowed to thee, O mighty God of Jacob, and the commandments which thou hast given unto him, and that he hath sincerely strove to do thy will.---Have mercy, O Lord, upon his wife and children, that they may be exalted in thy presence, and preserved by thy fostering hand.---Have mercy upon all their immediate connexions, that their prejudices may be broken up, and swept away as with a flood, that they may be converted and redeemed with Israel and know that thou art God. Remember, O, Lord, the presidents, even all the presidents of thy church, that thy right hand may exalt them with all their families, and their immediate connexions, that their names may be perpetuated and had in everlasting remembrance from generation to generation.

					Remember all thy church, O Lord, with all their families, and all their immediate connexions, with all their sick and afflicted ones, with all the poor and meek of the earth, that the kingdom which thou hast set up without hands, may become a great mountain and fill the whole earth, that thy church may come forth out of the wilderness of darkness, and shine forth fair as the moon, clear as the sun, and terrible as an army with banners, and be addorned as a bride for that day when thou shalt unveil the heavens, and cause the mountains to flow down at thy presence, and the valleys to be exalted, the rough places made smooth, that thy glory may fill the earth.

					That when the trump shall sound for the dead, we shall be caught up in the cloud to meet thee, that we may ever be with the Lord, that our garments may be pure, that we may be clothed upon with robes of righteousness, with palms in our hands, and crowns of glory upon our heads, and reap eternal joy for all our sufferings. O Lord, God Almighty. hear us in these our petitions, and answer us from heaven, thy holy habitation, where thou sittest enthroned, with glory, honor, power, majesty, might, dominion, truth, justice, judgement, mercy and an infinity of fulness, from everlasting to everlasting.

					O hear, O hear, O hear us, O Lord, and answer these petitions, and accept the dedication of this house, unto thee, the work of our hands, which we have built unto thy name; and also this church to put upon it thy name. And help us by the power of thy Spirit, that we may mingle our voices with those bright shining seraphs, around thy throne with acclamations of praise, singing hosanna to God and the Lamb: and let these thine annointed ones be clothed with salvation, and thy saints shout aloud for joy. \textsc{Amen and Amen}.
					
		% -----
		\subsubsection{Minutes and Prayer of Dedication, 27 March 1836 [D\&C 109]}
		% -----
			\label{EMS-1836-JS-27-MAR-1836-MIN}
			% \label{EMS-1830-JS-DAY-3LETTERMONTH-YEAR}
			
			\begin{description}
					\item[Print Sources]
				\end{description}

					\begin{tabular}{p{0.5in}p{1in}p{0.5in}p{1in}}
					JSP	 	& D5:188--209		& JSR 		& --- \\
					DC	 	& Section 109		& HC 		& 2:410--428 \\
					HDDC 	& 1440--1457		& TCHC 		& 2:399--414 \\
					\end{tabular}
					% HDDC = woodford's DC diss; JSR = Marquardt's JS Revelations; TCHC = Vogel HC
				
				\begin{description}
					\item[Online Access] \href{https://www.josephsmithpapers.org/paper-summary/minutes-and-prayer-of-dedication-27-march-1836-dc-109/1}{Here}
					% \item[Analysis] \hyperref[XX]{Here}
				\end{description}
				
				\begin{description}
					\item[Background]
				\end{description}
				
					% It might also be useful to give a two sentence context of the period: e.g., "This speech was given in Nauvoo to a small congregation one month prior to the destruction of the Nauvoo Expositor. These and other tensions may have been on the prophet's mind when he gave the speech."
				
				\begin{description}
					\item[Original Source]
				\end{description}
				
					% brief description of original source material, with reference to JSP or somewhere else for more info
					% Background material: it might be helpful to have a note that says whether a given document was presented as a speech, was written in a journal, etc.
					% Also a comment here about how reliable the source might be, or what shape it is in, if there are large omissions or parts that are hard to understand, and so on.
				
				\begin{description}
					\item[Text]
				\end{description}
				
						\hypertarget{EMS-1836-JS-27-MAR-1836-MIN-a}%
					...President
						\marginnote{a}%
					J. Smith jr. then rose, and after a few preliminary remarks, presented the several Presidents of the church, then present, to the several quorums respectively, and then to the church as being equal with himself, acknowledging them to be Prophets and Seers. The vote was unanimous in the affirmative in every instance.--- Each of the different quorums was presented in its turn to all the rest, and then to the church, and received and acknowledged by all the rest, in their several stations without a manifest dissenting sentiment.
					
					President J. Smith jr. then addressed the congregation in a manner calculated to instruct the understanding, rather than please the ear, and at or about the close of his remarks, he prophesied to all, that inasmuch as they would uphold these men in their several stations, alluding to the different quo[r]ums in the church, the Lord would bless them; yea, in the name of Christ, the blessings of Heaven shall be yours. And when the Lord's annointed go forth to proclaim the word, bearing testimony to this generation, if they receive it, they shall be blessed, but if not, the judgments of God will follow close upon them, until \textit{that} city or \textit{that} house, that rejects them, shall be left desolate...
					
						\hypertarget{EMS-1836-JS-27-MAR-1836-MIN-b}%
					...President
						\marginnote{b}%
					Smith then asked the several quorums separately and then the congregation, if they accepted the prayer. The vote was, in every instance, unanimous in the affirmative.
					
					The Eucharist was administered.--- D. C. [Don Carlos] Smith blessed the bread and wine and they were distributed by several Elders present, to the church.
					
					President J. Smith jr. then arose and bore record of his mission...
					
					...President S. Rigdon then made a few appropriate closing remarks; and a short prayer which was ended with loud acclamations of Hosanna! Hosanna! Hosanna to God and the Lamb, Amen, Amen and Amen! Three times. Elder B[righam] Young, one of the Twelve, gave a short address in tongues; Elder D[avid] W. Patten interpreted and gave a short exhortation in tongues himself; after which, President J. Smith jr. blessed the congregation in the name of the Lord, and at a little past four P. M. the whole exercise closed and the congregation dispersed...
					
		% -----
		\subsubsection{Record, 29 March 1836}
		% -----
			\label{EMS-1836-JS-29-MAR-1836-REC}
			% \label{EMS-1830-JS-DAY-3LETTERMONTH-YEAR}
			
			\begin{description}
					\item[Print Sources]
				\end{description}

					\begin{tabular}{p{0.5in}p{1in}p{0.5in}p{1in}}
					JSP	 	& ---				& JSR 		& --- \\
					DC	 	& ---				& HC 		& --- \\
					HDDC 	& ---				& TCHC 		& --- \\
					\end{tabular}
					% HDDC = woodford's DC diss; JSR = Marquardt's JS Revelations; TCHC = Vogel HC
				
				\begin{description}
					\item[Online Access] \href{https://www.josephsmithpapers.org/paper-summary/journal-1835-1836/188}{Here}
					% \item[Analysis] \hyperref[XX]{Here}
				\end{description}
				
				\begin{description}
					\item[Background]
				\end{description}
				
					% It might also be useful to give a two sentence context of the period: e.g., "This speech was given in Nauvoo to a small congregation one month prior to the destruction of the Nauvoo Expositor. These and other tensions may have been on the prophet's mind when he gave the speech."
				
				\begin{description}
					\item[Original Source]
				\end{description}
				
					% brief description of original source material, with reference to JSP or somewhere else for more info
					% Background material: it might be helpful to have a note that says whether a given document was presented as a speech, was written in a journal, etc.
					% Also a comment here about how reliable the source might be, or what shape it is in, if there are large omissions or parts that are hard to understand, and so on.
				
				\begin{description}
					\item[Text]
				\end{description}
				
						\hypertarget{EMS-1836-JS-29-MAR-1836-REC-a}%
					...\sout{After}
						\marginnote{a}%
					\sout{we} \sout{dismissed} \sout{made some arangements for our meeting on the morrow; attended to my domestick concirns, nothing very special transpired}

					\sout{At evening I met with the presidency in the Temple of the Lord and the Lord commanded us to tarry and santify ourselves by washing our feet}

						\hypertarget{EMS-1836-JS-29-MAR-1836-REC-b}%
					At
						\marginnote{b}%
					11 oclock A. M. Presidents Joseph Smith jun Frederick G. Williams, Sidney Rigdon, Hyrum Smith, and Oliver Cowdery met in the most holy place in the Lords house and sought for a revelation from Him, to teach us concerning our going to Zion, and other important matter after uniting in prayer, the voice of the Spirit was that we should come into this place three times, and also call the other presidents, the two Bishops and their councils (each to stand in his place) and fast through the day and also the night and that during this, if we would humble ourselves, we should receive further communication from Him.

					After this word was received, we immediately sent for the other brethren who came. The presidency proceeded to ordain George Boosinger to the high priesthood and annoint him.

					This was in consequence of his having administered unto us in temporal things in our distress. And also because he left the place just previous to the dedication of the Lords house to bring us the temporal means previously named.

						\hypertarget{EMS-1836-JS-29-MAR-1836-REC-c}%
					Soon
						\marginnote{c}%
					after this, the word of the Lord came to us through Presdt. J. Smith jun that those who had entered the holy place must not leave the house untill morning but send for such things as were necessary, and also, that during our stay we must cleans ourt feet and partake of the sacrament that we might be made holy before Him, and thereby be qualified to officiate in our calling upon the morrow in washing the feet of the Elders.

					Accordingly we proceeded and cleansed our faces and our feet, and then proceeded to wash each others feet.--- president S. Rigdon first washed presdt J. Smith jun and then in turn was washed by him--- after which president Rigdon washed presdt J. Smith Se\uline{n}. and Hyrum Smith <prsdt> J. Smith jun washed presdt F. G Williams, and then pres. Hyrum Smith washed president David Whitmer's feet and president Oliver Cowdery's, then pres D. Whitmer washed pres. W[illiam] W. Phelps feet and in turn pres Phelps washed pres John Whitmers feet.

					The Bishops and their councils were then washed: After which we partook of the bread and wine. The Holy S[p]irit rested down upon us and we continued in the Lords house all night prophesying and giving glory to God

		% -----
		\subsubsection{Revelation, 29 March 1836--A}
		% -----
			\label{EMS-1836-JS-29-MAR-1836-A}
			% \label{EMS-1830-JS-DAY-3LETTERMONTH-YEAR}
			
			\begin{description}
					\item[Print Sources]
				\end{description}

					\begin{tabular}{p{0.5in}p{1in}p{0.5in}p{1in}}
					JSP	 	& J1:212		& JSR 		& --- \\
					DC	 	& ---			& HC 		& 2:429 \\
					HDDC 	& ---			& TCHC 		& 2:415 \\
					\end{tabular}
					% HDDC = woodford's DC diss; JSR = Marquardt's JS Revelations; TCHC = Vogel HC
				
				\begin{description}
					\item[Online Access] \href{https://www.josephsmithpapers.org/paper-summary/revelation-29-march-1836-a/1}{Here}
					% \item[Analysis] \hyperref[XX]{Here}
				\end{description}
				
				\begin{description}
					\item[Background]
				\end{description}
				
					% It might also be useful to give a two sentence context of the period: e.g., "This speech was given in Nauvoo to a small congregation one month prior to the destruction of the Nauvoo Expositor. These and other tensions may have been on the prophet's mind when he gave the speech."
				
				\begin{description}
					\item[Original Source]
				\end{description}
				
					% brief description of original source material, with reference to JSP or somewhere else for more info
					% Background material: it might be helpful to have a note that says whether a given document was presented as a speech, was written in a journal, etc.
					% Also a comment here about how reliable the source might be, or what shape it is in, if there are large omissions or parts that are hard to understand, and so on.
				
				\begin{description}
					\item[Text]
				\end{description}
				
					the voice of the Spirit was that we should come into this place three times, and also call the other presidents, the two Bishops and their councils (each to stand in his place) and fast through the day and also the night and that during this, if we would humble ourselves, we should receive further communication from Him.
					
		% -----
		\subsubsection{Revelation, 29 March 1836--B}
		% -----
			\label{EMS-1836-JS-29-MAR-1836-B}
			% \label{EMS-1830-JS-DAY-3LETTERMONTH-YEAR}
			
			\begin{description}
					\item[Print Sources]
				\end{description}

					\begin{tabular}{p{0.5in}p{1in}p{0.5in}p{1in}}
					JSP	 	& J1:212		& JSR 		& --- \\
					DC	 	& ---			& HC 		& 2:430 \\
					HDDC 	& ---			& TCHC 		& 2:415 \\
					\end{tabular}
					% HDDC = woodford's DC diss; JSR = Marquardt's JS Revelations; TCHC = Vogel HC
				
				\begin{description}
					\item[Online Access] \href{https://www.josephsmithpapers.org/paper-summary/revelation-29-march-1836-b/1}{Here}
					% \item[Analysis] \hyperref[XX]{Here}
				\end{description}
				
				\begin{description}
					\item[Background]
				\end{description}
				
					% It might also be useful to give a two sentence context of the period: e.g., "This speech was given in Nauvoo to a small congregation one month prior to the destruction of the Nauvoo Expositor. These and other tensions may have been on the prophet's mind when he gave the speech."
				
				\begin{description}
					\item[Original Source]
				\end{description}
				
					% brief description of original source material, with reference to JSP or somewhere else for more info
					% Background material: it might be helpful to have a note that says whether a given document was presented as a speech, was written in a journal, etc.
					% Also a comment here about how reliable the source might be, or what shape it is in, if there are large omissions or parts that are hard to understand, and so on.
				
				\begin{description}
					\item[Text]
				\end{description}
				
					the word of the Lord came to us through Presdt. J. Smith jun that those who had entered the holy place must not leave the house untill morning but send for such things as were necessary, and also, that during our stay we must cleans ourt feet and partake of the sacrament that we might be made holy before Him, and thereby be qualified to officiate in our calling upon the morrow in washing the feet of the Elders.
					
		% -----
		\subsubsection{Certificate from Joshua Seixas, 30 March 1836}
		% -----
			\label{EMS-1836-JS-30-MAR-1836}
			% \label{EMS-1830-JS-DAY-3LETTERMONTH-YEAR}
			
			\begin{description}
					\item[Print Sources]
				\end{description}

					\begin{tabular}{p{0.5in}p{1in}p{0.5in}p{1in}}
					JSP	 	& D5:214--216		& JSR 		& --- \\
					DC	 	& ---				& HC 		& --- \\
					HDDC 	& ---				& TCHC 		& --- \\
					\end{tabular}
					% HDDC = woodford's DC diss; JSR = Marquardt's JS Revelations; TCHC = Vogel HC
				
				\begin{description}
					\item[Online Access] \href{https://www.josephsmithpapers.org/paper-summary/certificate-from-joshua-seixas-30-march-1836/1}{Here}
					% \item[Analysis] \hyperref[XX]{Here}
				\end{description}
				
				\begin{description}
					\item[Background]
				\end{description}
				
					% It might also be useful to give a two sentence context of the period: e.g., "This speech was given in Nauvoo to a small congregation one month prior to the destruction of the Nauvoo Expositor. These and other tensions may have been on the prophet's mind when he gave the speech."
				
				\begin{description}
					\item[Original Source]
				\end{description}
				
					% brief description of original source material, with reference to JSP or somewhere else for more info
					% Background material: it might be helpful to have a note that says whether a given document was presented as a speech, was written in a journal, etc.
					% Also a comment here about how reliable the source might be, or what shape it is in, if there are large omissions or parts that are hard to understand, and so on.
				
				\begin{description}
					\item[Text]
				\end{description}
				
					Mr Joseph Smith Jun\supers{r} has attended a full course of Hebrew lessons under my tuition; \& has been indefatigable in acquiring the principles of the sacred language of the Old Testament Scriptures in their original tongue. He has so far accomplished a knowledge of it, that he is able to translate to my entire satisfaction; \& by prosecuting the study he will be able to become a proficient in Hebrew. I take this opportunity of thanking him for his industry, \& his marked kindness towards me
					
					\begin{tightright}
					J[oshua] Seixas
					\end{tightright}
					
					Kirtland Ohio March 30\supers{th} 1836 [\textit{1/3 page blank}]
					
					<J. Seixas Recommendation>
					
					<Joseph Smith Junr.>
					
		% -----
		\subsubsection{Minutes, 30 March 1836}
		% -----
			\label{EMS-1836-JS-30-MAR-1836-MIN}
			% \label{EMS-1830-JS-DAY-3LETTERMONTH-YEAR}
			
			\begin{description}
					\item[Print Sources]
				\end{description}

					\begin{tabular}{p{0.5in}p{1in}p{0.5in}p{1in}}
					JSP	 	& D5:216--221; J1:213--216		& JSR 		& --- \\
					DC	 	& ---				& HC 		& 2:430--433 \\
					HDDC 	& ---				& TCHC 		& 2:416--418 \\
					TPJS 	& 109--110			& 	 		&  \\
					\end{tabular}
					% HDDC = woodford's DC diss; JSR = Marquardt's JS Revelations; TCHC = Vogel HC
				
				\begin{description}
					\item[Online Access] \href{https://www.josephsmithpapers.org/paper-summary/minutes-30-march-1836/1}{Here}
					% \item[Analysis] \hyperref[XX]{Here}
				\end{description}
				
				\begin{description}
					\item[Background]
				\end{description}
				
					% It might also be useful to give a two sentence context of the period: e.g., "This speech was given in Nauvoo to a small congregation one month prior to the destruction of the Nauvoo Expositor. These and other tensions may have been on the prophet's mind when he gave the speech."
				
				\begin{description}
					\item[Original Source]
				\end{description}
				
					% brief description of original source material, with reference to JSP or somewhere else for more info
					% Background material: it might be helpful to have a note that says whether a given document was presented as a speech, was written in a journal, etc.
					% Also a comment here about how reliable the source might be, or what shape it is in, if there are large omissions or parts that are hard to understand, and so on.
				
				\begin{description}
					\item[Text]
				\end{description}
				
						\hypertarget{EMS-1836-JS-30-MAR-1836-MIN-a}%
					Wedensday
						\marginnote{a}%
					morning 8 o clock March 30\supers{th} 1836 According to appointment the presidency, the 12, the seventies, the high \sout{councils} councils, the Bishops and their entire quorums, the Elders, and all the official members in this stake of Zion amounting to about 300 met in the temple of the Lord to attend to the ordinance of washing of feet, I ascended the pulpit and remarked to the congregation that we had passed through many trials and afflictions since the organization of this church and that this is a year of Jubilee to us and a time of rejoicing, and that it was expedient for us to prepare bread and wine sufficient to make our hearts glad, as we should not probably leave this house until morning; to this end we should call on the brethren to make a contribution, the stewards passed round and took up a liberal contribution and messengers were dispatched for bread and wine;
						\hypertarget{EMS-1836-JS-30-MAR-1836-MIN-b}%
					tubs
						\marginnote{b}%
					water and towels were prepared <and> I called the house to order, and the presidency proceeded to wash the feet of the 12 pronouncing many prophecy's and blessings upon them in the name of the Lord Jesus, the brethren began to prophesy upon each others heads, and cursings upon the enimies of Christ who inhabit Jackson county Missouri[.] continued prophesying and blessing and sealing them with Hosanna and Amen until nearly 7 o clock P.M. the bread <\& wine> was then brought in, and I observed that we had fasted all the day; and lest we faint; as the Saviour did so shall we do on this occasion, we shall bless the bread and give it to the 12 and they to the multitude, after which we shall bless the wine and do likewise; while waiting \sout{for} \sout{the} \sout{wine} I made the following remarks, that the time that we were required to tarry in Kirtland to be endued would be fulfilled in a few days, and then the Elders would go forth and each must stand for himself, that it was not necessary for them to be sent out two by two as in former times;
						\hypertarget{EMS-1836-JS-30-MAR-1836-MIN-c}%
					but
						\marginnote{c}%
					to go in all meekness in sobriety and preach Jesus Christ \& him crucified not to contend with others on the account of their faith or systems of religion but pursue a steady course, this I delivered by way of commandment, and all that observe them not will pull down persecution upon \sout{your} <thier> heads, while those who do shall always be filled with the Holy Ghost, this I pronounced as a prophesy, sealed with a Hosanna \& amen. Also that the seventies are not called to serve tables or preside over churches to settle difficulties, but to preach the gospel and build them up, and set others who do not belong to these quorums to preside over them who are high priests--- the twelve also are not to serve tables, but to bear the keys of the kingdom to all nations, and unlock them and call upon the seventies to follow after them and assist them. The 12 are at liberty to go wheresoever they will and if one shall say, I wish to go to such a place let all the rest say Amen.
					
						\hypertarget{EMS-1836-JS-30-MAR-1836-MIN-d}%
					The
						\marginnote{d}%
					seventies are at liberty to go to Zion if they please or go wheresoever they will and preach the gospel and let the redemtion of Zion be our object, and strive to affect it by sending up all the strength of the Lords house whereever we find them, and I want to enter into the following covenant, that if any more of our brethren are slain or driven from their lands in Missouri by the mob that we will give ourselves no rest until we are avenged of our enimies to the uttermost, this covenant was sealed unaminously by a hosanna and Amen.--- I then observed to the \sout{quorums}--- quorum<s> that I had now completed the\sout{ir} organization of the church and we had passed through all the necessary ceremonies, that I had given them all the instruction they needed and that they now were at liberty after obtaining their lisences%
						\footnote{%
							\label{EMS-1836-JS-30-MAR-1836-MIN-licenses}%
							Note that many licenses were issued on this day and the following days.
							}
					to go forth and build up the kingdom of God, and that it was expedient for me and the presidency to retire, having spent the night previous in waiting upon the Lord in his temple, and having to attend another dedication on the morrow, or conclude the one commenced on the last sabbath for the benifit of those of my brethren and sisters who could not get into the house on the former occasion but that it was expedient for the brethren to tarry all night and worship before the Lord in his house I left the meeting in the charge of the 12 and retired at about 9 o clock in the evening;
						\hypertarget{EMS-1836-JS-30-MAR-1836-MIN-e}%
					the
						\marginnote{e}%
					brethren continued exhorting, prophesying and speaking in tongues until 5 o clock in the morning--- the Saviour made his appearance to some, while angels minestered unto others, and it was a penticost and enduement indeed, long to be remembered for the sound shall go forth from this place into all the world, and the occurrences of this day shall be handed down upon the pages of sacred history to all generations, as the day of Pentecost, so shall this day be numbered and celebrated as a year of Jubilee and time of rejoicing to the saints of the most high God.
					
		% -----
		\subsubsection{Visions, 3 April 1836 [D\&C 110]}
		% -----
			\label{EMS-1836-JS-3-APR-1836}
			% \label{EMS-1830-JS-DAY-3LETTERMONTH-YEAR}
			
			\begin{description}
					\item[Print Sources]
				\end{description}

					\begin{tabular}{p{0.5in}p{1in}p{0.5in}p{1in}}
					JSP	 	& D5:224--229; J1:219--222		& JSR 		& 279--280 \\
					DC	 	& Section 110		& HC 		& 2:435--436 \\
					HDDC 	& 1458--1466		& TCHC 		& 2:420 \\
					\end{tabular}
					% HDDC = woodford's DC diss; JSR = Marquardt's JS Revelations; TCHC = Vogel HC
				
				\begin{description}
					\item[Online Access] \href{https://www.josephsmithpapers.org/paper-summary/visions-3-april-1836-dc-110/1}{Here}
					% \item[Analysis] \hyperref[DC110]{Here}
				\end{description}
				
				\begin{description}
					\item[Background]
				\end{description}
				
					% It might also be useful to give a two sentence context of the period: e.g., "This speech was given in Nauvoo to a small congregation one month prior to the destruction of the Nauvoo Expositor. These and other tensions may have been on the prophet's mind when he gave the speech."
					
					Need background here---include the first paragraph of the J1:219 entry
				
				\begin{description}
					\item[Original Source]
				\end{description}
				
					% brief description of original source material, with reference to JSP or somewhere else for more info
					% Background material: it might be helpful to have a note that says whether a given document was presented as a speech, was written in a journal, etc.
					% Also a comment here about how reliable the source might be, or what shape it is in, if there are large omissions or parts that are hard to understand, and so on.
				
				\begin{description}
					\item[Text]
				\end{description}
				
					The vail was taken from their minds and the eyes of their understandings were opened. They saw the Lord standing upon the breast work of the pulpit before them. and under his feet was a paved work of pure gold, in color like amber: his eyes were as a flame of fire; the hair of his head was like the pure snow, his countenance shone above the brightness of the sun, and his voice was as the sound of the rushing of great waters, even the Voice of Jehovah, saying, I am the first and the last. I am he who liveth. I am he who was slain. I am your Advocate with the Father. Behold your sins are forgiven you. You are clean before me, therefore, lift up your heads and rejoice, let the hearts of your brethren rejoice and let the hearts of all my \sout{brethren} <people> rejoice, who have with their might, built this house to my name. For behold I have accepted this house and my name shall be here; and I will manifest myself to my people, in mercy, in this House, Yea I will appear unto my servants and speak unto them with mine own voice, if my people will keep my commandments and do not pollute this Holy House. Yea the hearts of thousands and tens of thousands shall greatly rejoice in consequence of the blessings which shall be poured out, and the endowment with which my servants have already been endowed and shall hereafter be endowed in this House. and the fame of this House shall spread to foreign lands, and this is the beginning of the blessing, which shall be poured out upon the heads of my people. even so amen. After this vision closed, the Heavens were again opened unto them and Moses appeared before them and committed unto them the Keys of the gathering of Israel from the four parts of the Eearth and the leading of the ten tribes from the Land of the North. After this Elias appeared and committed the dispensation of the gospel of Abraham, saying, that in them and their seed all generations after them should be blessed. After this vision had closed, another great and glorious vision burts [burst] upon them, for Elijah, the Prophet, who was taken to Heaven without tasting death, also stood before them, and said, behold the time has fully come which was spoken of by the mouth of Malachi, testifying, that he should be sent before the great and dreadful day of the Lord come, to turn the hearts of the Fathers to the children, and the children to the fathers, lest the whole earth be smitten with a curse. Therefore, the Keys of this dispensation are committed into your hands, and by this ye may know that the great and the dreadful day of the Lord is near, even at the doors
					
		% -----
		\subsubsection{Letter to Oliver Cowdery, circa 9 April 1836}
		% -----
			\label{EMS-1836-JS-9-APR-1836}
			% \label{EMS-1830-JS-DAY-3LETTERMONTH-YEAR}
			
			\begin{description}
					\item[Print Sources]
				\end{description}

					\begin{tabular}{p{0.5in}p{1in}p{0.5in}p{1in}}
					JSP	 	& D5:231--243		& JSR 		& --- \\
					DC	 	& ---				& HC 		& 2:436--440 \\
					HDDC 	& ---				& TCHC 		& 2:421--425 \\
					\end{tabular}
					% HDDC = woodford's DC diss; JSR = Marquardt's JS Revelations; TCHC = Vogel HC
				
				\begin{description}
					\item[Online Access] \href{https://www.josephsmithpapers.org/paper-summary/letter-to-oliver-cowdery-circa-9-april-1836/1}{Here}
					% \item[Analysis] \hyperref[XX]{Here}
				\end{description}
				
				\begin{description}
					\item[Background]
				\end{description}
				
					% It might also be useful to give a two sentence context of the period: e.g., "This speech was given in Nauvoo to a small congregation one month prior to the destruction of the Nauvoo Expositor. These and other tensions may have been on the prophet's mind when he gave the speech."
				
				\begin{description}
					\item[Original Source]
				\end{description}
				
					% brief description of original source material, with reference to JSP or somewhere else for more info
					% Background material: it might be helpful to have a note that says whether a given document was presented as a speech, was written in a journal, etc.
					% Also a comment here about how reliable the source might be, or what shape it is in, if there are large omissions or parts that are hard to understand, and so on.
				
				\begin{description}
					\item[Text]
				\end{description}
				
					\begin{tightcenter}
							\hypertarget{EMS-1836-JS-9-APR-1836-a}%
						\phantom{ }\textit{For}
							\marginnote{a}%
						\textit{the Messenger and Advocate.}
							
					\end{tightcenter}
					
					\textsc{Brother O}[liver] \textsc{Cowdery:}

					\textit{Dear Sir}---This place having recently been visited by a gentleman who advocated the principles or doctrines of those who are called abolitionists; if you deem the following reflections of any service, or think they will have a tendency to correct the opinions of the southern public, relative to the views and sentiments I believe, as an individual, and am able to say, from personal knowledge, are the feelings of others, you are at liberty to give them publicity in the columns of the Advocate. I am prompted to this course in consequence, in one respect, of many elders having gone into the Southern States, besides, there now being many in that country who have already embraced the fulness of the gospel, as revealed through the book of Mormon,---having learned, by experience, that the enemy of truth does not slumber, nor cease his exertions to bias the minds of communities against the servants of the Lord, by stiring up the indignation of men upon all matters of importance or interest.

						\hypertarget{EMS-1836-JS-9-APR-1836-b}%
					Thinking,
						\marginnote{b}%
					perhaps, that the sound might go out, that ``an abolitionist'' had held forth several times to this community, and that the public feeling was not aroused to create mobs or disturbances, leaving the impression that all he said was concurred in, and received as gospel and the word of salvation. I am happy to say, that no violence or breach of the public peace was attempted, so far from this, that all except a very few, attended to their own avocations and left the gentleman to hold forth his own arguments to nearly naked walls.

						\hypertarget{EMS-1836-JS-9-APR-1836-c}%
					I
						\marginnote{c}%
					am aware, that many who profess to preach the gospel, complain against their brethren of the same faith, who reside in the south, and are ready to withdraw the hand of fellowship because they will not renounce the principle of slavery and raise their voice against every thing of the kind. This must be a tender point, and one which should call forth the candid reflection of all men, and especially before they advance in an opposition calculated to lay waste the fair States of the South, and set loose, upon the world a community of people who might peradventure, overrun our country and violate the most sacred principles of human society,---chastity and virtue.

					No one will pretend to say, that the people of the free states are as capable of knowing the evils of slavery as those who hold them. If slavery is an evil, who, could we expect, would first learn it? Would the people of the free states, or would the slave states? All must readily admit, that th[e] latter would first learn this fact. If the fact was learned first by those immediately concerned, who would be more capable than they of prescribing a remedy?

						\hypertarget{EMS-1836-JS-9-APR-1836-d}%
					And
						\marginnote{d}%
					besides, are not those who hold slaves, persons of ability, discernment and candor? Do they not expect to give an account at the bar of God for their conduct in this life? It may, no doubt, with propriety be said, that many who hold slaves live without the fear of God before their eyes, and, the same may be said of many in the free states. Then who is to be the judge in this matter?

					So long, then, as those of the free states are not interested in the freedom of the slaves, any other than upon the mere principles of equal rights and of the gospel, and are ready to admit that there are men of piety who reside in the South, who are immediately concerned, and until \textit{they} complain, and call for assistance, why not cease their clamor, and no further urge the slave to acts of murder, and the master to vigorous discipline, rendering both miserable, and unprepared to pursue that course which might otherwise lead them both to better their condition? I do not believe that the people of the North have any more right to say that the South \textit{shall not} hold slaves, than the South have to say the North \textit{shall}.

						\hypertarget{EMS-1836-JS-9-APR-1836-e}%
					And
						\marginnote{e}%
					further, what benefit will it ever be to the slave for persons to run over the free states, and excite indignation against their masters in the minds of thousands and tens of thousands who understand nothing relative to their circumstances or conditions? I mean particularly those who have never travelled in the South, and scarcely seen a negro in all their life. How any community can ever be excited with the chatter of such persons---boys and others who are too indolent to obtain their living by honest industry, and are incapable of pursuing any occupation of a professional nature, is unaccountable to me. And when I see persons in the free states signing documents against slavery, it is no less, in my mind, than an array of influence, and a declaration of hostilities against the people of the South! What can divide our Union sooner, God only knows!

						\hypertarget{EMS-1836-JS-9-APR-1836-f}%
					After
						\marginnote{f}%
					having expressed myself so freely upon this subject, I do not doubt but those who have been forward in raising their voice against the South, will cry out against me as being uncharitable, unfeeling and unkind---wholly unacquainted with the gospel of Christ. It is my privilege then, to name certain passages from the bible, and examine the teachings of the ancients upon this matter, as the fact is uncontrovertable, that the first mention we have of slavery is found in the holy bible, pronounced by a man who was perfect in his generation and walked with God. And so far from that prediction's being averse from the mind of God it remains as a lasting monument of the decree of Jehovah, to the shame and confusion of all who have cried out against the South, in consequence of their holding the sons of Ham in servitude!

					``And he said cursed \textit{be} Canaan; a servant of servants shall he be unto his brethren. And he said, Blessed \textit{be} the Lord God of Shem; and Canaan shall be his servant.--- God shall enlarge Japheth, and he shall dwell in the tents of Shem; and Canaan shall be his servant.''---Gen, 8:25, 26, 27.

						\hypertarget{EMS-1836-JS-9-APR-1836-g}%
					Trace
						\marginnote{g}%
					the history of the world from this notable event down to this day, and you will find the fulfilment of this singular prophecy. What could have been the design of the Almighty in this wonderful occurrence is not for me to say; but I can say, that the curse is not yet taken off the sons of Canaan, neither will be until it is affected by as great power as caused it to come; and the people who interfere the least with the decrees and purposes of God in this matter, will come under the least condemnation before him; and those who are determined to pursue a course which shows an opposition and a feverish restlessness against the designs of the Lord, will learn, when perhaps it is too late for their own good, that God can do his own work without the aid of those who are not dictated by his counsel.

						\hypertarget{EMS-1836-JS-9-APR-1836-h}%
					I
						\marginnote{h}%
					must not pass over a notice of the history of Abraham, of whom so much is spoken in the scriptures. If we can credit the account, God conversed with him from time to time, and directed him in the way he should walk, saying, ``I am the Almighty God: walk before me and be thou perfect.'' Paul says that the gospel was preached to this man. And it is further said, that he had sheep and oxen, men-servants and maid-servants, \&c. From this I conclude, that if the principle had been an evil one, in the midst of the communications made to this holy man, he would have been instructed differently. And if he was instructed against holding men-servants and maid-servants, he never ceased to do it; consequently must have incurred the displeasure of the Lord and thereby lost his blessings---which was not the fact.

						\hypertarget{EMS-1836-JS-9-APR-1836-i}%
					Some
						\marginnote{i}%
					may urge, that the names, man-servant and maid-servant, only mean hired persons who were at liberty to leave their masters or employers at any time. But we can easily settle this point by turning to the history of Abraham's descendants, when governed by a law given from the mouth of the Lord himself. I know that when an Israelite had been brought into servitude in consequence of debt, or otherwise, at the seventh year he went from the task of his former master or employer; but to no other people or nation was this granted in the law to Israel. And if, after a man had served six years, he did not wish to be free, then the master was to bring him unto the judges, boar his ear with an awl, and that man was ``to serve him forever.'' The conclusion I draw from this, is that this people were led and governed by revelation and if such a law was wrong God only is to be blamed, and abolitionists are not responsible.

						\hypertarget{EMS-1836-JS-9-APR-1836-j}%
					Now,
						\marginnote{j}%
					before proceeding any farther, I wish to ask one or two questions:---Were the apostles men of God, and did they preach the gospel? I have no doubt but those who believe the bible will admit these facts, and that they also knew the mind and will of God concerning what they wrote to the churches which they were instrumental in building up.

					This being admitted, the matter can be put to rest without much argument, if we look at a few items in the New Testament. Paul says:

					``Servants, be obedient to them that are \textit{your} masters according to the flesh, with fear and trembling, in singleness of your heart, as unto Christ: Not with eye service, as men-pleasers: but as the servants of Christ, doing the will of God from the heart: With good will doing service, as to the Lord, and not to men. Knowing that whatsoever good thing any man doeth, the same shall he receive of the Lord, whether \textit{he be} bond or free. And, ye masters, do the same things unto them, forbearing threatening: knowing that your Master also is in heaven; neither is there respect of persons with him.'' Eph. 6:5, 6, 7, 8, 9.

						\hypertarget{EMS-1836-JS-9-APR-1836-k}%
					Here
						\marginnote{k}%
					is a lesson which might be profitable for all to learn, and the principle upon which the church was anciently governed, is so plainly set forth, that an eye of truth might see and understand. Here, certainly are represented the master and servant; and so far from instructions to the servant to leave his master, he is commanded to be in obedience, as unto the Lord: the master in turn is required to treat them with kindness before God, understanding at the same time that he is to give an account.--- The hand of fellowship is not withdrawn from him in consequence of having servants.

						\hypertarget{EMS-1836-JS-9-APR-1836-l}%
					The
						\marginnote{l}%
					same wri[t]er, in his first epistle to Timothy, the sixth chapter, and the five first verses, says:

					``Let as many servants as are under the yoke count their own masters worthy of all honor, that the name of God and \textit{his} doctrine be not blasphemed. And they that have believing masters, let them not despise \textit{them}, because they are brethren: but rather do \textit{them} service, because they are faithful and beloved, partakers of the benefit These things teach and exhort. If any man teach otherwise, and consent not to wholesome words, \textit{even} the words of our Lord Jesus Christ, and to the doctrine which is according to godliness: he is proud, knowing nothing but doting about questions and strifes of words, whereof cometh envy, strife, railings, evil surmisings, Perverse disputings of men of corrupt minds, and destitute of the truth, supposing that gain is godliness: from such withdraw thyself.''

					This is so perfectly plain, that I see no need of comment. The scripture stands for itself, and I believe that these men were better qualified to teach the will of God, than all the abolitionists in the world.

						\hypertarget{EMS-1836-JS-9-APR-1836-m}%
					Before
						\marginnote{m}%
					closing this communication, I beg leave to drop a word to the travelling elders: You know, brethren, that great responsibility rests upon you, and that you are accountable to God for all you teach the world. In my opinion, you will do well to search the book of Covenants, in which you will see the belief of the church concerning masters and servants. All men are to be taught to repent; but we have no right to interfere with slaves contrary to the mind and will of their masters. In fact, it would be much better and more prudent, not to preach at all to slaves, until after their masters are converted: and then, teach the master to use them with kindness, remembering that they are accountable to God, and that servants are bound to serve their masters, with singleness of heart, without murmuring. I do, most sincerely hope, that no one who is authorized from this church to preach the gospel, will so far depart from the scripture as to be found stirring up strife and sedition against our brethren of the South. Having spoken frankly and freely, I leave all in the hands of God, who will direct all things for his glory and the accomplishment of his work.

					Praying that God may spare you to do much good in this life, I subscribe myself your brother in the Lord.

					\begin{tightright}
					JOSEPH SMITH, jr.
					\end{tightright}
					
		% -----
		\subsubsection{Letter to John Thornton and Others, 25 July 1836}
		% -----
			\label{EMS-1836-JS-25-JUL-1836}
			% \label{EMS-1830-JS-DAY-3LETTERMONTH-YEAR}
			
			\begin{description}
					\item[Print Sources]
				\end{description}

					\begin{tabular}{p{0.5in}p{1in}p{0.5in}p{1in}}
					JSP	 	& D5:268--271		& JSR 		& --- \\
					DC	 	& ---				& HC 		& 2:456--461 \\
					HDDC 	& ---				& TCHC 		& 2:442--447 \\
					\end{tabular}
					% HDDC = woodford's DC diss; JSR = Marquardt's JS Revelations; TCHC = Vogel HC
				
				\begin{description}
					\item[Online Access] \href{https://www.josephsmithpapers.org/paper-summary/letter-to-john-thornton-and-others-25-july-1836/1}{Here}
					% \item[Analysis] \hyperref[XX]{Here}
				\end{description}
				
				\begin{description}
					\item[Background]
				\end{description}
				
					% It might also be useful to give a two sentence context of the period: e.g., "This speech was given in Nauvoo to a small congregation one month prior to the destruction of the Nauvoo Expositor. These and other tensions may have been on the prophet's mind when he gave the speech."
				
				\begin{description}
					\item[Original Source]
				\end{description}
				
					% brief description of original source material, with reference to JSP or somewhere else for more info
					% Background material: it might be helpful to have a note that says whether a given document was presented as a speech, was written in a journal, etc.
					% Also a comment here about how reliable the source might be, or what shape it is in, if there are large omissions or parts that are hard to understand, and so on.
				
				\begin{description}
					\item[Text]
				\end{description}
				
						\hypertarget{EMS-1836-JS-25-JUL-1836-a}%
					...Another
						\marginnote{a}%
					charge brought against our friends is that of being dangerous in societies ``where slavery is tolerated and practiced.'' Without occupying time here, we refer you to the April (1836) No. of the ``Latter Day Saints' Messenger and Advocate,'' printed at this place, a copy of which we forward to each of you. From the length of time which has transpired since its publication, you can easily see, that it was put forth for no other reason than to correct the public mind generally, without a reference or expectation of an excitement of the nature of the one now in your country. Why we refer you to this publication, particularly, is because many of our friends who are now at the west, were in this place when this paper made its appearance, and from personal observation gave it their decided approbation, and expressed those sentiments to be their own, in the fullest particular...
						
						\hypertarget{EMS-1836-JS-25-JUL-1836-b}%
					...If
						\marginnote{b}%
					a people, a community, or a society, can accumulate wealth, increase in worldly fortune, improve in science and arts, rise to eminence in the eyes of the public, surmount these difficulties, so much as to bid defiance to poverty and wretchedness, it must be a new creation, a race of beings super-human. But in all their poverty and want, we have yet to learn, for the first time, that our friends are not industrious, and temperate, and wherein they have not always been the \textit{last} to retaliate or resent an injury, and the \textit{first} to overlook and forgive. We do not urge that there are not exceptions to be found: all communities, all societies and associations, are cumbered with disorderly and less virtuous members---members who violate in a greater or less degree the principles of the same. But this can be no just criterion by which to judge a whole society. And further still, where a people are laboring under constant fear of being dispossessed, very little inducement is held out to excite them to be industrious...
					
						\hypertarget{EMS-1836-JS-25-JUL-1836-c}%
					...Though
						\marginnote{c}%
					our religious principles are before the world, ready for the investigation of all men, yet we are aware that the sole foundation of all the persecution against our friends, has arisen in consequence of the calumnies and misconstructions, without foundation in truth, or righteousness, in common with all other religious societies, at their first commencement; and should Providence order that we rise not as others before us, to respectability and esteem, but be trodden down by the ruthless hand of extermination, \textit{posterity} will do us the justice, when our persecutors are equally low in the dust, with ourselves, to hand down to succeeding generations, the virtuous acts and forbearance of a people, who sacrificed their reputation for their religion, and their earthly fortunes and happiness, to preserve peace, and save this land from being further drenched in blood.
					
						\hypertarget{EMS-1836-JS-25-JUL-1836-d}%
					We
						\marginnote{d}%
					have no doubt but your very seasonable mediation, in the time of so great an excitement, will accomplish your most sanguine desire, in preventing further disorder; and we hope, gentlemen, that while you reflect upon the fact, that the citizens of Clay county are \textit{urgent} for our friends to leave you, that you will also bear in mind, that by their complying with your request to leave, is surrendering some of the dearest rights and first, among those inherent principles, guaranteed in the constitution of our country; and that human nature can be driven to a certain extent, when it will yield no farther. Therefore, while our friends \textit{suffer} so much, and forego so many sacred rights, we sincerely hope, and we have every reason to expect it, that a suitable forbearance may be shown by the people of Clay, which if done, the cloud that has been obscuring your horizon, will disperse, and you be left to enjoy peace, harmony and prosperity...
					
		% -----
		\subsubsection{Revelation, 6 August 1836 [D\&C 111]}
		% -----
			\label{EMS-1836-JS-6-AUG-1836}
			% \label{EMS-1830-JS-DAY-3LETTERMONTH-YEAR}
			
			\begin{description}
					\item[Print Sources]
				\end{description}

					\begin{tabular}{p{0.5in}p{1in}p{0.5in}p{1in}}
					JSP	 	& D5:271--278		& JSR 		& 281 \\
					DC	 	& Section 111		& HC 		& 2:465--466 \\
					HDDC 	& 1467--1474		& TCHC 		& 2:450--451 \\
					\end{tabular}
					% HDDC = woodford's DC diss; JSR = Marquardt's JS Revelations; TCHC = Vogel HC
				
				\begin{description}
					\item[Online Access] \href{https://www.josephsmithpapers.org/paper-summary/revelation-6-august-1836-dc-111/1}{Here}
					% \item[Analysis] \hyperref[XX]{Here}
				\end{description}
				
				\begin{description}
					\item[Background]
				\end{description}
				
					% It might also be useful to give a two sentence context of the period: e.g., "This speech was given in Nauvoo to a small congregation one month prior to the destruction of the Nauvoo Expositor. These and other tensions may have been on the prophet's mind when he gave the speech."
				
				\begin{description}
					\item[Original Source]
				\end{description}
				
					% brief description of original source material, with reference to JSP or somewhere else for more info
					% Background material: it might be helpful to have a note that says whether a given document was presented as a speech, was written in a journal, etc.
					% Also a comment here about how reliable the source might be, or what shape it is in, if there are large omissions or parts that are hard to understand, and so on.
				
				\begin{description}
					\item[Text]
				\end{description}
				
					\begin{tightcenter}
					A revelation
					
					Salem (Mss.) [Massachusetts] August 6, 1836.
					\end{tightcenter}
					
					I the Lord your God am not displeased with your coming this Journey, notwithstandi[n]g your follies. I have much treasure in this city for you, for the benefit of Zion; and many people in this city whom I will gather out in due time for the benefit of Zion, through your instrumentality: Therefore it is expedient that you should form acquaintance with men in this city, as you shall be lead, and as it shall be be given you. And it shall come to pass, in due time, that I will give this city into your hands, that you shall have power over it, insomuch that they shall not discover your secret parts; and its wealth, pertaining to gold and silver, shall be yours. Concern not yourselves about your debts, for I will give you power to pay them. Concern not yourselves about Zion, for I will deal merciful with her. Tarry in this place and in the regions round about, and the place where it is my will that you should tarry, for the main, shall be signalized unto you by the peace and power of \sout{the} my Spirit, that shall flow unto you. This place you may obtain by hire, \&c . . . And inquire diligently concerning the more ancient inhabitants and founders of this city, for there are more treasures than one for you, in this city: Therefore be ye as wise as serpents and yet without sin, and I will order all things for your good as fast as ye are able to receive them. Amen.

% -----------------------------------------------------------
\pagebreak{}
% -----------------------------------------------------------